\documentclass[11pt]{article}

\usepackage[letterpaper,margin=0.82in,headheight=15pt]{geometry}
\usepackage[T1]{fontenc}
\usepackage{lmodern}
\usepackage{microtype}
\usepackage{amsmath,amssymb,amsthm,bm,mathtools}
\usepackage{booktabs,longtable,tabularx,array,multirow}
\usepackage{enumitem}
\usepackage{xcolor}
\usepackage{fancyhdr}
\usepackage{graphicx}
\usepackage{tikz-cd}
\usepackage{hyperref}
\usepackage[nameinlink,noabbrev]{cleveref}

\definecolor{navy}{HTML}{17324D}
\definecolor{teal}{HTML}{147D92}
\definecolor{orange}{HTML}{D47A27}
\definecolor{softblue}{HTML}{EEF4F7}
\definecolor{softgray}{HTML}{F3F5F6}
\definecolor{darkgray}{HTML}{46545E}
\definecolor{softorange}{HTML}{FFF4E8}

\hypersetup{
  colorlinks=true,
  linkcolor=navy,
  citecolor=teal,
  urlcolor=teal,
  pdftitle={Volume IV: Matched Spin-1 Nuclear Light-Front Dynamics and Deuteron GTMDs},
  pdfauthor={DeuteronWigner project}
}

\pagestyle{fancy}
\fancyhf{}
\fancyhead[L]{\small\color{darkgray}DeuteronWigner formalism - Volume IV}
\fancyhead[R]{\small\color{darkgray}Matched spin-1 nuclear dynamics}
\fancyfoot[L]{\small\color{darkgray}29 July 2026}
\fancyfoot[R]{\small\color{darkgray}\thepage}
\renewcommand{\headrulewidth}{0.4pt}

\setlist[itemize]{leftmargin=1.5em,itemsep=2pt,topsep=3pt}
\setlist[enumerate]{leftmargin=1.7em,itemsep=2pt,topsep=3pt}
\setlength{\parindent}{0pt}
\setlength{\parskip}{5pt}
\setlength{\emergencystretch}{2em}
\renewcommand{\arraystretch}{1.16}
\allowdisplaybreaks

\newtheorem{definition}{Definition}[section]
\newtheorem{axiom}[definition]{Axiom}
\newtheorem{principle}[definition]{Design principle}
\newtheorem{requirement}[definition]{Requirement}
\newtheorem{proposition}[definition]{Proposition}
\newtheorem{theorem}[definition]{Theorem}
\newtheorem{lemma}[definition]{Lemma}
\newtheorem{corollary}[definition]{Corollary}
\newtheorem{remark}[definition]{Remark}

\crefname{definition}{definition}{definitions}
\crefname{axiom}{axiom}{axioms}
\crefname{principle}{principle}{principles}
\crefname{requirement}{requirement}{requirements}
\crefname{proposition}{proposition}{propositions}
\crefname{theorem}{theorem}{theorems}
\crefname{lemma}{lemma}{lemmas}
\crefname{corollary}{corollary}{corollaries}
\crefname{remark}{remark}{remarks}

\newcommand{\kT}{\bm{k}_{T}}
\newcommand{\ellT}{\bm{\ell}_{T}}
\newcommand{\pT}{\bm{p}_{T}}
\newcommand{\qT}{\bm{q}_{T}}
\newcommand{\PT}{\bm{P}_{T}}
\newcommand{\DT}{\bm{\Delta}_{T}}
\newcommand{\DNT}{\bm{\Delta}_{N T}}
\newcommand{\bD}{\bm{b}_{\Delta}}
\newcommand{\bM}{\bm{b}_{\mathrm{TMD}}}
\newcommand{\RT}{\bm{R}_{T}}
\newcommand{\dd}{\mathrm{d}}
\newcommand{\Tr}{\operatorname{Tr}}
\newcommand{\tr}{\operatorname{tr}}
\newcommand{\Span}{\operatorname{span}}
\newcommand{\End}{\operatorname{End}}
\newcommand{\Hom}{\operatorname{Hom}}
\newcommand{\Id}{\operatorname{Id}}
\newcommand{\Ker}{\operatorname{Ker}}
\newcommand{\Ran}{\operatorname{Ran}}
\newcommand{\rank}{\operatorname{rank}}
\newcommand{\diag}{\operatorname{diag}}
\newcommand{\Hil}{\mathcal H}
\newcommand{\Phys}{\mathcal H_{\mathrm{phys}}}
\newcommand{\Alg}{\mathcal A}
\newcommand{\Rep}{\mathcal V}
\newcommand{\LF}{\mathrm{LF}}
\newcommand{\TMD}{\mathrm{TMD}}
\newcommand{\QCD}{\mathrm{QCD}}
\newcommand{\GTMD}{\mathrm{GTMD}}
\newcommand{\GPD}{\mathrm{GPD}}
\newcommand{\EFT}{\mathrm{EFT}}
\newcommand{\Amp}{\mathbf{Amp}}
\newcommand{\Dens}{\mathbf{Dens}}
\newcommand{\Match}{\mathbf{Match}}
\newcommand{\Red}{\mathbf{Red}}
\newcommand{\Proc}{\mathbf{Proc}}
\newcommand{\cO}{\mathcal O}
\newcommand{\cM}{\mathcal M}
\newcommand{\cR}{\mathcal R}
\newcommand{\cS}{\mathcal S}
\newcommand{\cV}{\mathcal V}
\newcommand{\cW}{\mathcal W}
\newcommand{\cE}{\mathcal E}
\newcommand{\cG}{\mathcal G}
\newcommand{\cK}{\mathcal K}
\newcommand{\cQ}{\mathcal Q}
\newcommand{\cP}{\mathcal P}
\newcommand{\cC}{\mathcal C}
\newcommand{\cZ}{\mathcal Z}
\newcommand{\cD}{\mathcal D}
\newcommand{\cN}{\mathcal N}
\newcommand{\cU}{\mathcal U}
\newcommand{\cF}{\mathcal F}
\newcommand{\cA}{\mathcal A}
\newcommand{\cB}{\mathcal B}
\newcommand{\cL}{\mathcal L}
\newcommand{\cJ}{\mathcal J}
\newcommand{\cI}{\mathcal I}
\newcommand{\eps}{\varepsilon}
\newcommand{\bbone}{\bm 1}
\newcommand{\abs}[1]{\left|#1\right|}
\newcommand{\norm}[1]{\left\lVert#1\right\rVert}
\newcommand{\inner}[2]{\left\langle #1\middle|#2\right\rangle}
\newcommand{\avg}[1]{\left\langle #1\right\rangle}
\newcommand{\comm}[2]{\left[#1,#2\right]}
\newcommand{\acom}[2]{\left\{#1,#2\right\}}
\newcommand{\crossT}[2]{\bigl(#1\!\times\!#2\bigr)_z}
\newcommand{\NN}{NN}
\newcommand{\NNpi}{NN\pi}
\newcommand{\DD}{\Delta\Delta}
\newcommand{\sixq}{6q}
\newcommand{\IA}{\mathrm{IA}}
\newcommand{\coh}{\mathrm{coh}}
\newcommand{\exch}{\mathrm{exch}}
\newcommand{\shortd}{\mathrm{short}}
\newcommand{\nuc}{\mathrm{nuc}}
\newcommand{\micro}{\mathrm{micro}}
\newcommand{\ren}{\mathrm{ren}}
\newcommand{\phys}{\mathrm{phys}}
\newcommand{\spec}{\mathrm{spec}}
\newcommand{\tagg}{\mathrm{tag}}
\newcommand{\off}{\mathrm{off}}
\newcommand{\ct}{\mathrm{ct}}
\newcommand{\soft}{\mathrm{soft}}
\newcommand{\rap}{\mathrm{rap}}
\newcommand{\UV}{\mathrm{UV}}
\newcommand{\IR}{\mathrm{IR}}
\newcommand{\even}{\mathrm{even}}
\newcommand{\odd}{\mathrm{odd}}

\newcolumntype{Y}{>{\raggedright\arraybackslash}X}
\newcolumntype{L}[1]{>{\raggedright\arraybackslash}p{#1}}

\title{\vspace{-1.0cm}
{\color{navy}\bfseries Volume IV: Matched Spin-1 Nuclear Light-Front Dynamics\\[2mm]
and Deuteron GTMDs}\\[4mm]
\large Nuclear Fock space, quark--gluon-to-hadronic operator matching,\\
coherent amplitudes, non-nucleonic sectors, and tensor-polarized reductions}
\author{DeuteronWigner project}
\date{29 July 2026}

\begin{document}
\maketitle

\begin{center}
\begin{minipage}{0.94\textwidth}
\colorbox{softblue}{\parbox{0.96\textwidth}{
\textbf{Status and purpose.}
Volume~I defined regulated light-front state spaces and microscopic dynamics.
Volume~II constructed common quark, antiquark, and gluon nucleon GTMD overlaps,
and Volume~III supplied link-resolved rescattering amplitudes.  The present
volume is the matched nuclear layer.  Its primary realizable route is a
hierarchy in which the quark--gluon nucleon theory is matched onto a
spin-resolved nuclear light-front effective theory with explicit
\(NN\), \(NN\pi\), \(\Delta\Delta\), six-quark, and controlled additional
sectors.  The deuteron GTMD is a matrix element in one normalized spin-1 state.
Impulse approximation, exchange currents, coherent shadowing, pion exchange,
and non-nucleonic mechanisms are derived reductions or operator blocks of
that state, not independently normalized curves.  The accepted fixed-scale
canonical model remains a regression and comparison boundary; no existing
phenomenological component is promoted to microscopic status by notation.
}}
\end{minipage}
\end{center}

\begin{abstract}
A genuinely predictive deuteron partonic model cannot be obtained by
convolving a microscopic nucleon correlator with an unrelated nuclear
response.  The nuclear Hamiltonian, state, current, and partonic operators
must belong to one matched low-resolution theory.  This volume constructs
that theory as the spin-1 continuation of the common nucleon GTMD formalism.
The primary route is hierarchical: a regulated quark--gluon nucleon model is
matched at a hadronic resolution onto a nuclear light-front model space with
nucleon, pion, Delta, short-distance, and supported compact sectors.  The same
matching transformation generates the nuclear Hamiltonian and the one-body,
two-body, mesonic, transition, coherent, and short-distance partonic
operators.  Unitarity and field-redefinition covariance show why a standalone
``off-shell nucleon correction'' is not observable: strength can be moved
between off-shell one-body terms and many-body operators while the complete
matrix element remains invariant.

We define the regulated spin-1 nuclear Hilbert space, the block mass operator,
one normalized deuteron Fock expansion, baryon/charge/plus-momentum ledgers,
and a directed nuclear truncation system.  The \(NN\) component is obtained
from a rotationally covariant rest-frame \(S\)- and \(D\)-wave amplitude by a
declared Jacobian and constituent light-front spin rotations.  At zero
skewness we derive the exact two-body recoil map: in a spectator-preserving
one-body matrix element the active nucleon receives the full external
transverse transfer, while the intrinsic deuteron wave-function arguments
shift by \((1-y)\bm\Delta_T/2\).  Two-body operators instead carry an explicit
transfer-partition variable.  The resulting off-forward spectral amplitude
is kept distinct from the partonic GTMD and from the TMD impact variable.

The full deuteron parent is
\(
 W_{a/D}^{[\Gamma,\gamma]}
 =\sum_{\alpha,\beta}
 \langle\Psi_D^\alpha|\widehat{\mathcal O}_{a,\alpha\beta}^{[\Gamma,\gamma]}
 |\Psi_D^\beta\rangle
\),
with resolved nuclear sectors, nucleon helicities, flavor, Wilson path,
gluon color class, transfer, transverse rank, matching scheme, and common
member identity.  The familiar light-front impulse convolution is derived as
the spectator-preserving one-body reduction of this expression and is shown
to share its TMD, GPD, PDF, current, and form-factor marginals.  Spin-1
irreducible projectors expose how nucleon vector polarization, \(S\)-\(D\)
interference, orbital harmonics, and many-body operators populate the
\(U,L,T,LL,LT,TT\) channels.  The \(b_1\) relation is retained as the first
inclusive tensor anchor, while elastic charge, magnetic, and quadrupole form
factors, the light-front angular condition, tagged spectator observables, and
momentum/current sum rules provide independent closure tests.

Coherent small-\(x\) shadowing is formulated as a helicity-resolved
multiple-scattering amplitude using the same deuteron state and diffractive
nucleon amplitudes; polarized and tensor shadowing cannot be copied from an
unpolarized ratio.  The \(NN\pi\) sector carries a genuine spin-resolved
three-body amplitude and includes pion-active, nucleon-active, and transition
operators with a subtraction that prevents double counting the nucleon's
internal pion cloud.  \(\Delta\Delta\), six-quark, hidden-color, and
short-range sectors require explicit amplitudes and operators before they may
enter a nonzero central prediction.  Completely positive maps are derived
only after a justified partial trace; coherent sector interference remains at
amplitude level.  Partonic Wilson transport and nuclear rescattering are kept
separate and connected by an explicit overlap subtraction whenever their
soft regions coincide.

The volume concludes with formal propositions, analytic benchmarks,
machine-readable interfaces, property-based tests, a staged C4 implementation
program, and acceptance criteria for export to the common matching,
evolution, and process layer of Volume~V.
\end{abstract}

\tableofcontents

\section{Role, scope, and inherited contracts}
\label{sec:role}

\subsection{The change of dynamical level}

Volumes~I--III constructed the nucleon-side object
\begin{equation}
 \mathfrak W_{a/N}
 =\left(
 W_{a/N,\lambda'\lambda}^{[J,\gamma]},
 \mathfrak o_{a/N},
 \mathfrak B_{a/N},
 \mathfrak m_N,
 \operatorname{Cov}_{\micro}
 \right),
\label{eq:nucleon-input}
\end{equation}
where \(N=p,n\), \(J\) denotes the quark Dirac or gluon Lorentz projection,
\(\gamma\) is the complete Wilson geometry, \(\mathfrak B\) is the phase and
soft-overlap budget, and \(\mathfrak m_N\) is the shared microscopic member.
The nuclear problem is not to dress a named scalar TMD.  It is to evaluate the
operator represented by \cref{eq:nucleon-input} in a normalized spin-1 state
whose Hamiltonian, currents, and many-body operators are matched to the same
hadronic resolution.

The required direction of inference is
\begin{equation}
 \left
 \{H_{\LF}^{qg},\,|\Psi_N\rangle,\,\cO_{a/N}^{qg}\right\}
 \xrightarrow{\;\cZ_{qg\to\nuc}\;}
 \left\{H_{\LF}^{\nuc},\,\widehat\cO_a^{\nuc}\right\}
 \xrightarrow{\;\mathrm{solve}\;}
 |\Psi_D\rangle
 \xrightarrow{\;\mathrm{matrix\ element}\;}
 W_{a/D}.
\label{eq:source-chain}
\end{equation}
This is a matched hierarchy, not a claim that the first implementation solves
the six-quark continuum deuteron directly.  A direct quark--gluon deuteron
solution remains a later ultraviolet completion against which the hierarchy
can be tested.

\begin{principle}[One matched nuclear theory]
\label{prin:one-nuclear-theory}
The nuclear Hamiltonian, state, local currents, and partonic operators must
share a resolution, regulator, unitary convention, and matching record.  A
high-quality deuteron wave function and a high-quality operator from a
different theory do not form a controlled prediction merely because each is
accurate separately.
\end{principle}

\subsection{Primary domain}

The first complete implementation uses
\begin{equation}
 \xi_D=-\frac{\Delta_D^+}{2P_D^+}=0,
 \qquad
 \Delta_D^+=0,
 \qquad
 t_D=-\DT^2,
\label{eq:xi-zero-domain}
\end{equation}
with a spin-1 deuteron target and quark, antiquark, and gluon operators at
leading twist.  The domain is sufficient for:
\begin{itemize}
\item the deuteron TMD limit and partonic Wigner transform;
\item zero-skewness GPDs and elastic/current moments;
\item the complete forward \(U,L,T,LL,LT,TT\) projection system;
\item inclusive \(b_1\), tensor PDFs, and tagged spectator structure;
\item a first coherent small-\(x\) amplitude and explicit non-nucleonic sectors.
\end{itemize}
Nonzero skewness, breakup GTMDs, transition GPDs, and complete polynomiality
are later extensions, but the operator and recoil types below are designed so
that they can be added without reinterpreting existing data.

\subsection{Inherited fail-closed contracts}

The following Volume~II and Volume~III rules remain mandatory:
\begin{enumerate}
\item The nuclear layer consumes nucleon helicity matrices, not only named
      scalar distributions.
\item The external deuteron transfer \(\DT\), active-nucleon transfer
      \(\DNT\), nuclear intrinsic momentum \(\pT\), partonic momentum \(\kT\),
      nuclear Wigner coordinate \(\RT\), partonic Wigner coordinate \(\bD\),
      and TMD impact parameter \(\bM\) are distinct types.
\item Proton/neutron, flavor, link, gluon ordered-link and \(f/d\) color,
      transverse-rank, scheme, and member identities remain resolved.
\item The partonic staple \(U_\gamma^{\rm parton}\) is not identified with
      coherent nuclear propagation \(\cU_{\nuc}\).
\item Off-diagonal nuclear-sector matrix elements are allowed only when the
      matched operator changes the corresponding sector.
\item A completely positive map may be introduced only after the relevant
      coherent degrees of freedom are no longer resolved.
\end{enumerate}

\subsection{Explicit nonclaims}

This volume does not claim:
\begin{itemize}
\item a continuum six-quark solution of the deuteron;
\item that realistic \(S/D\) wave functions alone determine partonic
      two-body operators;
\item that an ``off-shell nucleon PDF'' is a separately measurable object;
\item that unpolarized diffractive shadowing fixes polarized or tensor
      shadowing;
\item that a scalar pion light-cone distribution is sufficient for all
      tensor TMDs;
\item that every nuclear correction is an incoherent quantum channel;
\item that hidden color or \(\Delta\Delta\) probability alone specifies a
      flavor-, spin-, color-, and transfer-resolved GTMD;
\item that completion of this formalism by itself establishes numerical
      convergence or predictive accuracy.
\end{itemize}

\section{Resolution hierarchy and matched nuclear model space}
\label{sec:matching-hierarchy}

\subsection{Separated physical scales}

The deuteron is weakly bound, with a binding momentum parametrically much
smaller than the pion mass and hadronic internal scale.  The hierarchy
\begin{equation}
 \gamma_D\ll m_\pi\lesssim \Lambda_{\nuc}\ll \Lambda_{qg}
\label{eq:scale-hierarchy}
\end{equation}
motivates a low-resolution nuclear theory with nucleons, pions, Delta
isobars, and contact operators, while nucleon quark--gluon substructure is
retained through matched one-hadron operators.  The separation is not exact:
pion and soft-gluon regions can occur on both sides of the matching surface.
Those overlap regions require explicit subtraction rather than verbal
assignment.

A nuclear resolution datum is
\begin{equation}
 \mathfrak r_D=
 \left(
 \Lambda_{\UV}^{D},\Lambda_{\IR}^{D},
 \mathcal B_D,\mathcal F_D,
 \nu_H,\nu_O,
 \cR_D,\cS_D
 \right),
\label{eq:nuclear-resolution-datum}
\end{equation}
where \(\mathcal B_D\) is the basis or grid, \(\mathcal F_D\) the retained
nuclear Fock sectors, \(\nu_H\) and \(\nu_O\) the Hamiltonian and operator
orders, and \(\cR_D,\cS_D\) the regulator and renormalization schemes.

\subsection{Model-space projection and effective operators}

Let \(P_D\) project onto the retained nuclear model space and \(Q_D=1-P_D\).
For a full Hamiltonian \(H\), the exact projected resolvent defines the
energy-dependent effective Hamiltonian
\begin{equation}
 H_{\rm eff}(E)
 =P_DHP_D
 +P_DHQ_D\frac{1}{E-Q_DHQ_D}Q_DHP_D.
\label{eq:feshbach-hamiltonian}
\end{equation}
The same elimination changes every operator.  For a full operator \(\cO\),
one convenient wave-operator representation is
\begin{equation}
 \cO_{\rm eff}(E_f,E_i)
 =P_D\Omega^\dagger(E_f)\,\cO\,\Omega(E_i)P_D,
\label{eq:effective-operator}
\end{equation}
with \(\Omega\) the exact map from model-space vectors to full-space states.
Expanding \cref{eq:effective-operator} produces one-body, exchange,
transition, and contact terms even when the original operator is one-body.

\begin{requirement}[Hamiltonian--operator co-matching]
\label{req:ham-operator-comatching}
Whenever a sector or momentum region is eliminated from the Hamiltonian, the
induced contribution to every current and partonic operator used in the same
calculation must be retained to the declared order.  Eliminating the sector
from \(H\) while keeping a bare one-body \(\cO\) is not a matched
approximation.
\end{requirement}

\subsection{Unitary and field-redefinition covariance}

Let \(U_D\) be a unitary transformation within the retained nuclear space.
Then
\begin{equation}
 H_D' =U_DH_DU_D^\dagger,
 \qquad
 \cO_D'=U_D\cO_DU_D^\dagger,
 \qquad
 |\Psi_D'\rangle=U_D|\Psi_D\rangle
\label{eq:unitary-covariance}
\end{equation}
give identical complete matrix elements,
\begin{equation}
 \langle\Psi_D'|\cO_D'|\Psi_D'\rangle
 =\langle\Psi_D|\cO_D|\Psi_D\rangle.
\label{eq:unitary-matrix-invariance}
\end{equation}
Field redefinitions are the Lagrangian counterpart of this statement and can
move strength between off-shell two-body interactions, many-body forces, and
many-body currents without changing observables
\cite{FurnstahlHammerTirfessa2001,Krebs2020}.

\begin{proposition}[No standalone off-shell observable]
\label{prop:no-offshell-observable}
Suppose a one-body nucleon operator is changed by a term proportional to the
nucleon equation of motion,
\begin{equation}
 \delta\cO_N=(p^2-M_N^2)X.
\label{eq:eom-offshell-change}
\end{equation}
A field redefinition can remove \(\delta\cO_N\) and generate a compensating
many-body operator \(\delta\cO_D^{(2+)}\).  Therefore the complete deuteron
matrix element is invariant, while the separately quoted ``off-shell
correction'' is representation dependent.
\end{proposition}

\begin{proof}
Apply the equivalence transformation that eliminates the redundant
operator \((p^2-M_N^2)X\).  The transformed action contains induced
many-body terms.  Equation~\eqref{eq:unitary-matrix-invariance} then implies
that only the sum of one- and many-body contributions is invariant.  Omitting
the induced term creates an artificial representation dependence.
\end{proof}

\subsection{Primary and direct routes}

The primary implementation route is
\begin{equation}
 \text{quark--gluon nucleon theory}
 \xrightarrow{\;\cZ_{qg\to\mathrm{had}}\;}
 \left(H_D^{N,\pi,\Delta,\ldots},\widehat{\cO}_a^D\right)
 \xrightarrow{\;\mathrm{solve}\;}
 |\Psi_D\rangle
 \xrightarrow{\;\mathrm{evaluate}\;}
 W_{a/D}.
\label{eq:hierarchical-route}
\end{equation}
A later direct route solves for a deuteron in a quark--gluon Hilbert space
with six-quark, gluon, sea, cluster, and hadronic asymptotic sectors.  The
nuclear EFT must then emerge as a matched low-resolution reduction of that
same state.  Results from unrelated direct and hierarchical models may be
compared, but may not be spliced without a matching map.

\section{Nuclear light-front kinematics and typed recoil}
\label{sec:kinematics}

\subsection{External deuteron momentum fibers}

We inherit
\begin{equation}
 v^\pm=\frac{v^0\pm v^3}{\sqrt2},
 \qquad
 v\cdot w=v^+w^-+v^-w^+-\bm v_T\cdot\bm w_T.
\label{eq:lf-convention}
\end{equation}
At zero skewness choose the symmetric frame
\begin{align}
 P_i&=P-\frac{\Delta}{2},&
 P_f&=P+\frac{\Delta}{2},
 \notag\\
 P_{iT}&=-\frac{\DT}{2},&
 P_{fT}&=+\frac{\DT}{2},&
 \Delta^+&=0.
\label{eq:deuteron-symmetric-frame}
\end{align}
The deuteron operator acts between distinct momentum fibers,
\begin{equation}
 \widehat\cO_a(P_f,P_i):\Hil_{D,P_i}\longrightarrow\Hil_{D,P_f}.
\label{eq:deuteron-operator-fibers}
\end{equation}

The canonical partonic fraction is
\begin{equation}
 x_D=\frac{k^+}{P_D^+},\qquad 0<x_D<1.
\label{eq:xD-definition}
\end{equation}
The project production coordinate is \(x_N=2x_D\).  Number-preserving density
adapters therefore obey
\begin{equation}
 f^{[x_N]}(x_N)=\frac12 f^{[x_D]}\!\left(\frac{x_N}{2}\right),
\label{eq:x-adapter}
\end{equation}
so that the support, number, and momentum ledgers are not changed by a
coordinate relabeling.

\subsection{Two-body intrinsic variables and measure}

For the \(NN\) sector, let nucleon 1 carry fraction \(y\) and average intrinsic
momentum \(\pT\), while nucleon 2 carries \(1-y\) and \(-\pT\).  The free
invariant mass is
\begin{equation}
 M_0^2(y,\pT)
 =\frac{M_1^2+\pT^2}{y}
 +\frac{M_2^2+\pT^2}{1-y}.
\label{eq:two-body-free-mass}
\end{equation}
For equal masses,
\begin{equation}
 k_z=\left(y-\frac12\right)M_0,
 \qquad
 \frac{\partial k_z}{\partial y}
 =\frac{M_0}{4y(1-y)}.
\label{eq:kz-jacobian}
\end{equation}
We use the invariant two-body measure
\begin{equation}
 [\dd\Gamma_2]
 =\frac{\dd y\,\dd^2\pT}{2(2\pi)^3y(1-y)}.
\label{eq:two-body-measure}
\end{equation}
Any alternative normalization is represented by a unitary adapter and must
transform the wave function, current, spectral amplitude, and convolution
Jacobian together.

\subsection{Exact spectator-preserving recoil map}

Let \(\pT\) denote the average intrinsic momentum of the active nucleon.  In
the one-body impulse block the physical spectator momentum is unchanged.
The intrinsic arguments are then
\begin{align}
 \bm p_{1T}^{\rm in}
 &=\pT-\frac{1-y}{2}\DT,
 &
 \bm p_{1T}^{\rm out}
 &=\pT+\frac{1-y}{2}\DT,
\label{eq:active-intrinsic-recoil}\\
 \bm p_{2T}^{\rm in}
 &=-\pT+\frac{1-y}{2}\DT,
 &
 \bm p_{2T}^{\rm out}
 &=-\pT-\frac{1-y}{2}\DT.
\label{eq:spectator-intrinsic-recoil}
\end{align}
Combining these with the external hadron momenta gives
\begin{align}
 p_{1T}^{\rm phys,in}&=\pT-\frac12\DT,&
 p_{1T}^{\rm phys,out}&=\pT+\frac12\DT,
\notag\\
 p_{2T}^{\rm phys,in}&=-\pT,&
 p_{2T}^{\rm phys,out}&=-\pT.
\label{eq:physical-one-body-recoil}
\end{align}
Thus the active nucleon receives the full physical transfer,
\begin{equation}
 \DNT=\DT
 \qquad\text{for the spectator-preserving one-body block},
\label{eq:active-full-transfer}
\end{equation}
while the deuteron wave-function arguments shift by \((1-y)\DT/2\).
Conflating these two statements is a common recoil error.

\begin{proposition}[Two-body recoil closure]
\label{prop:recoil-closure}
Equations~\eqref{eq:active-intrinsic-recoil}--\eqref{eq:spectator-intrinsic-recoil}
preserve intrinsic transverse closure separately in the initial and final
states and leave the spectator physical momentum unchanged.
\end{proposition}

\begin{proof}
The initial intrinsic momenta sum to
\(
 \pT-(1-y)\DT/2-\pT+(1-y)\DT/2=0
\), and similarly for the final state.  Inserting the external share
\((1-y)P_{i,f,T}\) into the spectator momentum gives \(-\pT\) on both sides.
The active physical momentum therefore differs by \(\DT\).
\end{proof}

\subsection{Transfer partition for many-body operators}

A two-body operator can transfer momentum to both constituents.  Introduce a
partition variable \(\bm\delta_T\) such that
\begin{equation}
 \bm\Delta_{1T}=\bm\delta_T,
 \qquad
 \bm\Delta_{2T}=\DT-\bm\delta_T.
\label{eq:two-body-transfer-partition}
\end{equation}
The matched kernel determines the distribution in \(\bm\delta_T\).  Setting
\(\bm\delta_T=\DT\) is the one-body spectator limit; setting
\(\bm\delta_T=y\DT\) or any other partition without deriving it from the
operator is forbidden.

A typed nuclear recoil object contains
\begin{equation}
 \mathfrak r_{\nuc}
 =\left(
 P_i,P_f,
 y,\pT,
 \DT,\DNT,
 \bm\delta_T,
 \text{active slot},
 \text{spectator policy},
 \text{Jacobian convention}
 \right).
\label{eq:nuclear-recoil-type}
\end{equation}

\begin{requirement}[Single recoil authority]
\label{req:single-recoil-authority}
Wave functions, spectral amplitudes, one-body GTMD convolutions, two-body
operators, tagged observables, and form-factor moments must consume the same
versioned nuclear recoil object.  No projector or benchmark may reimplement
its shift factors locally.
\end{requirement}

\subsection{Partonic momentum inside the active nucleon}

Let \(z=x_D/y\) be the parton fraction of the active nucleon.  If \(\kT\) is
the parton's average transverse momentum relative to the deuteron, the
nucleon-intrinsic parton momentum is
\begin{equation}
 \ellT=\kT-z\pT.
\label{eq:parton-nucleon-transverse-map}
\end{equation}
The mapping is a boost relation, not a model for transverse broadening.  Its
forward and off-forward use must share the same convention so that TMD, GPD,
PDF, and current reductions commute.

\section{Regulated spin-1 nuclear Hilbert space and dynamics}
\label{sec:nuclear-hilbert}

\subsection{Sector decomposition}

At nuclear resolution \(\mathfrak r_D\), the physical model space is
\begin{equation}
 \Hil_D^{(\mathfrak r_D)}
 =\Hil_{\NN}
 \oplus\Hil_{\NNpi}
 \oplus\Hil_{\DD}
 \oplus\Hil_{\sixq}
 \oplus\Hil_{\rm other}.
\label{eq:nuclear-hilbert-sum}
\end{equation}
At the charge-symmetric reference point, every retained component belongs to
the deuteron channel
\begin{equation}
 B=2,
 \qquad Q=+1,
 \qquad I=0,
 \qquad J^P=1^+.
\label{eq:deuteron-quantum-numbers}
\end{equation}
Baryon number, charge, and total angular momentum/parity are physical sector
labels.  The value \(I=0\) defines the charge-symmetric reference space;
controlled \(m_d-m_u\) and electromagnetic operators may admix small
\(I\neq0\) components, which must be retained as explicitly tagged symmetry-
breaking amplitudes rather than folded silently into the isoscalar state.  A
sector label includes constituent species, helicities, isospin couplings,
orbital basis, color coupling where relevant, regulator, and matching origin.

The absence of an \(N\Delta\) sector from the central \(I=0\) model space is
not a numerical simplification: \(\tfrac12\otimes\tfrac32=1\oplus2\) does not
contain total isospin zero.  Delta degrees of freedom first enter through
allowed transition operators, \(NN\pi\)-like sectors, or \(\Delta\Delta\).
An \(N\Delta\) component could appear only in a separately declared
isospin-breaking extension and is not part of the charge-symmetric central
state.

\subsection{Block mass operator}

The nuclear light-front mass operator is
\begin{equation}
 \cM_D^2
 =\cM_{0,D}^2+V_D+\delta\cM_D^2,
 \qquad
 \cM_D^2=(\cM_D^2)^\dagger,
\label{eq:nuclear-mass-operator}
\end{equation}
with sector blocks
\begin{equation}
 \cM_{D,\alpha\beta}^2
 =P_\alpha\cM_D^2P_\beta,
 \qquad
 \alpha,\beta\in
 \{\NN,\NNpi,\DD,\sixq,\ldots\}.
\label{eq:nuclear-mass-blocks}
\end{equation}
The diagonal blocks contain sector kinetic energies and interactions; the
off-diagonal blocks generate pion emission/absorption, pair or compact-sector
mixing, and other retained transitions.  Every block has a Hermitian partner,
a regulator identity, and a matching/subtraction record.

The eigenproblem is
\begin{equation}
 \cM_D^2|D,P,\Lambda;\mathfrak m_D\rangle
 =M_D^2|D,P,\Lambda;\mathfrak m_D\rangle,
\label{eq:deuteron-eigenproblem}
\end{equation}
where \(\mathfrak m_D\) is the complete nuclear member, not merely a wave
function name.

\subsection{One normalized spin-1 state}

Write
\begin{equation}
 |D,\Lambda\rangle
 =|\Psi_{\NN}^{\Lambda}\rangle
 +|\Psi_{\NNpi}^{\Lambda}\rangle
 +|\Psi_{\DD}^{\Lambda}\rangle
 +|\Psi_{\sixq}^{\Lambda}\rangle
 +\cdots .
\label{eq:deuteron-fock-expansion}
\end{equation}
Orthogonality of distinct Fock sectors gives
\begin{equation}
 1=\langle D,\Lambda|D,\Lambda\rangle
 =Z_{\NN}+Z_{\NNpi}+Z_{\DD}+Z_{\sixq}+\cdots,
\qquad
 Z_\alpha=\norm{\Psi_\alpha^\Lambda}^2.
\label{eq:sector-normalization}
\end{equation}
The sector probabilities are resolution and scheme dependent; the total
normalization and conserved charges are physical constraints.

The plus-momentum ledger is
\begin{equation}
 1=
 \sum_\alpha
 \int[\dd\Gamma_\alpha]
 \left(\sum_{i\in\alpha}y_i\right)
 |\Psi_\alpha|^2,
\label{eq:nuclear-plus-ledger}
\end{equation}
and the baryon and electric-charge ledgers are
\begin{align}
 2&=
 \sum_\alpha\int[\dd\Gamma_\alpha]
 \left(\sum_{i\in\alpha}B_i\right)|\Psi_\alpha|^2,
\label{eq:nuclear-baryon-ledger}\\
 1&=
 \sum_\alpha\int[\dd\Gamma_\alpha]
 \left(\sum_{i\in\alpha}Q_i\right)|\Psi_\alpha|^2.
\label{eq:nuclear-charge-ledger}
\end{align}
The normalization convention must not be repaired separately in each
observable.

\subsection{Nuclear truncation tower}

A nuclear truncation label is
\begin{equation}
 r_D=
 (\Lambda_D,N_D,\mathcal F_D,
 \nu_H,\nu_J,\nu_{\cO},n_{\coh},n_{\rm grid}),
\label{eq:nuclear-truncation-label}
\end{equation}
where \(\nu_J\) and \(\nu_{\cO}\) are local-current and partonic-operator
orders and \(n_{\coh}\) is the multiple-scattering order.  Calculations at
different \(r_D\) are compared through matched observables,
\begin{equation}
 \delta_i(r_D',r_D)
 =\left\|
 R_{r_D'\to r_D}[\cO_i^{(r_D')}]
 -\cO_i^{(r_D)}
 \right\|,
\label{eq:nuclear-convergence-residual}
\end{equation}
not by demanding regulator-independent sector probabilities or bare radial
wave functions.

\begin{requirement}[Independent nuclear error axes]
\label{req:nuclear-error-axes}
Basis, regulator, EFT order, current order, partonic-operator order, Fock
content, coherent-scattering order, non-nucleonic content, and quadrature
errors remain separate.  Convergence in the \(NN\) wave function cannot be
used to hide missing two-body operators or missing sectors.
\end{requirement}

\section{The nucleonic deuteron component}
\label{sec:nn-component}

\subsection{Light-front two-nucleon state}

With the measure in \cref{eq:two-body-measure}, the nucleonic component is
\begin{equation}
 |\Psi_{\NN}^{\Lambda}(P)\rangle
 =\sum_{\lambda_p,\lambda_n}
 \int[\dd\Gamma_2]\,
 \Psi_{\Lambda;\lambda_p\lambda_n}^{\LF}(y,\pT)
 |p(y,\pT,\lambda_p)\,n(1-y,-\pT,\lambda_n)\rangle,
\label{eq:nn-state}
\end{equation}
with the isosinglet factor understood.  In a charge-symmetric convention,
exchange of the two nucleons obeys the required fermionic antisymmetry after
spin, isospin, and orbital factors are combined.

The \(NN\) normalization is
\begin{equation}
 \sum_{\lambda_p,\lambda_n}
 \int[\dd\Gamma_2]\,
 \abs{\Psi_{\Lambda;\lambda_p\lambda_n}^{\LF}(y,\pT)}^2
 =Z_{\NN}.
\label{eq:nn-normalization}
\end{equation}

\subsection{Rest-frame S- and D-wave amplitude}

A rotationally covariant rest-frame amplitude is
\begin{align}
 \Phi_{\Lambda;\sigma_p\sigma_n}(\bm k)
 ={}&\sum_{L=0,2}u_L(k)
 \sum_{m_L,m_S}
 \langle Lm_L,1m_S|1\Lambda\rangle
 Y_{Lm_L}(\widehat{\bm k})
 \notag\\
 &\times
 \left\langle\frac12\sigma_p,\frac12\sigma_n\middle|1m_S\right\rangle
 \chi_{I=0}.
\label{eq:instant-sd-wave}
\end{align}
Here \(u_0\) and \(u_2\) are the \(S\)- and \(D\)-wave radial components.
The probability \(P_D\) is model and resolution dependent; signed
\(S\)-\(D\) interference is physical in matrix elements and cannot be
reconstructed from \(P_D\) alone.

\subsection{Canonical-to-light-front spin map}

For a constituent of mass \(m\), fraction \(y\), and transverse momentum
\(\pT\), define the Melosh rotation
\begin{equation}
 R_M(y,\pT)
 =\frac{m+yM_0
 -i\bm\sigma\cdot(\widehat{\bm z}\times\pT)}
 {\sqrt{(m+yM_0)^2+\pT^2}}.
\label{eq:melosh-rotation}
\end{equation}
The light-front amplitude is
\begin{align}
 \Psi_{\Lambda;\lambda_p\lambda_n}^{\LF}(y,\pT)
 ={}&
 \sqrt{\frac{\partial k_z}{\partial y}}
 \sum_{\sigma_p,\sigma_n}
 D^{1/2}_{\lambda_p\sigma_p}[R_M(y,\pT)]
 \notag\\
 &\times
 D^{1/2}_{\lambda_n\sigma_n}[R_M(1-y,-\pT)]
 \Phi_{\Lambda;\sigma_p\sigma_n}(\bm k),
\label{eq:lf-wavefunction-map}
\end{align}
up to the declared normalization adapter.  This is a specific light-front
realization, not a theorem that every front-form construction has only two
radial components.  Explicitly covariant light-front dynamics can contain
additional invariant amplitudes, and their omission must be counted as a
truncation axis \cite{Carbonell1998}.

\begin{requirement}[Wave-function identity]
\label{req:wavefunction-identity}
A nuclear wave-function object records the rest-frame interaction, radial
normalization, canonical-to-light-front map, spin-rotation convention,
relativistic components retained, regulator, current family, and numerical
member.  The label ``AV18,'' ``CD-Bonn,'' or ``Norfolk'' alone is incomplete.
\end{requirement}

\subsection{Rotational covariance and the angular condition}

In front-form calculations, transverse rotations are dynamical.  A truncated
wave function and one-body current may therefore produce spurious dependence
on the light-front orientation.  For the plus current in a \(q^+=0\) frame,
four independent helicity matrix elements must reduce to three physical form
factors.  In one standard convention the angular-condition residual is
\begin{equation}
 \Delta_{\rm ang}(Q^2)
 =(1+2\eta)I_{11}^+
 +I_{1,-1}^+
 -\sqrt{8\eta}\,I_{10}^+
 -I_{00}^+,
 \qquad
 \eta=\frac{Q^2}{4M_D^2}.
\label{eq:angular-condition}
\end{equation}
A covariant complete calculation has \(\Delta_{\rm ang}=0\).  At finite
truncation, the residual is a diagnostic of missing rotational components,
exchange currents, pair terms, or zero modes, not a quantity to remove by
choosing whichever form-factor extraction prescription is most favorable
\cite{Carbonell1998,CookeMiller2002}.

\subsection{Controlled wave-function family}

The first numerical family may include AV18, CD-Bonn, and local chiral/Norfolk
interactions \cite{Wiringa1995,Machleidt2001,Piarulli2016}.  They are
alternative nuclear members.  They are not simultaneous amplitudes in one
state and are not added in quadrature as if they were independent statistical
replicas.  A central member must use currents and matching operators
consistent with its interaction; cross-family comparisons define a model or
EFT sensitivity axis.

\section{Spin-1 irreducible amplitudes and nuclear spectral geometry}
\label{sec:spectral-geometry}

\subsection{Irreducible target operators}

Let \(|1\Lambda\rangle\), \(\Lambda=+1,0,-1\), span the deuteron helicity
space.  A trace-orthonormal irreducible tensor basis is
\begin{equation}
 \left(T_M^{(K)}\right)_{\Lambda'\Lambda}
 =\sqrt{\frac{2K+1}{3}}
 \langle 1\Lambda,KM|1\Lambda'\rangle,
 \qquad
 K=0,1,2,
\label{eq:spin1-irrep-basis}
\end{equation}
with
\begin{equation}
 \Tr\left[
 (T_M^{(K)})^\dagger T_{M'}^{(K')}
 \right]
 =\delta_{KK'}\delta_{MM'}.
\label{eq:spin1-irrep-orthogonality}
\end{equation}
The scalar, vector, and quadrupole components of any deuteron matrix \(X\) are
\begin{equation}
 X_M^{(K)}=\Tr[(T_M^{(K)})^\dagger X].
\label{eq:spin1-irrep-coefficients}
\end{equation}
The real \(U,L,T,LL,LT,TT\) basis is obtained through a fixed invertible
adapter.  In particular,
\begin{equation}
 \delta_T X
 =X_{00}-\frac12(X_{++}+X_{--})
\label{eq:tensor-helicity-difference}
\end{equation}
is convention independent, while the named \(LL\) coefficient depends on the
chosen \(S_{LL}\) normalization.

\subsection{Off-forward nucleon spectral amplitude}

For an active nucleon \(N\) and a spectator nucleon \(s\), define the
spin-resolved off-forward spectral amplitude
\begin{align}
 \cS_{\Lambda'\Lambda;\lambda'\lambda}^{N/D}
 (y,\pT,\DT)
 ={}&
 \sum_{\lambda_s}
 \Psi_{\Lambda';\lambda'\lambda_s}^{\LF *}
 \left(y,\pT+\frac{1-y}{2}\DT\right)
 \notag\\
 &\times
 \Psi_{\Lambda;\lambda\lambda_s}^{\LF}
 \left(y,\pT-\frac{1-y}{2}\DT\right),
\label{eq:offforward-spectral-amplitude}
\end{align}
with the normalization factors implied by
\cref{eq:two-body-measure,eq:nn-normalization}.  This object is a matrix in
both deuteron and active-nucleon helicities.  It is not a scalar smearing
function.

Hermiticity gives
\begin{equation}
 \cS_{\Lambda'\Lambda;\lambda'\lambda}^{N/D}
 (y,\pT,\DT)
 =
 \left[
 \cS_{\Lambda\Lambda';\lambda\lambda'}^{N/D}
 (y,\pT,-\DT)
 \right]^*.
\label{eq:spectral-hermiticity}
\end{equation}
At \(\DT=0\),
\begin{equation}
 \rho_{\Lambda'\Lambda;\lambda'\lambda}^{N/D}(y,\pT)
 =\cS_{\Lambda'\Lambda;\lambda'\lambda}^{N/D}(y,\pT,0)
\label{eq:forward-spectral-density}
\end{equation}
is positive semidefinite in the combined deuteron--active-nucleon space when
constructed from the complete wave function.

\subsection{Nuclear Wigner transform}

The nuclear transverse Wigner amplitude is
\begin{equation}
 \cW_{\Lambda'\Lambda;\lambda'\lambda}^{N/D}
 (y,\pT,\RT)
 =\int\frac{\dd^2\DT}{(2\pi)^2}
 e^{-i\RT\cdot\DT}
 \cS_{\Lambda'\Lambda;\lambda'\lambda}^{N/D}(y,\pT,\DT).
\label{eq:nuclear-wigner-transform}
\end{equation}
It satisfies
\begin{equation}
 \int\dd^2\RT\,
 \cW^{N/D}(y,\pT,\RT)
 =\rho^{N/D}(y,\pT).
\label{eq:nuclear-wigner-marginal}
\end{equation}
The coordinate \(\RT\) is conjugate to nuclear constituent transfer; it is
not the partonic Wigner coordinate \(\bD\) and not the TMD separation
\(\bM\).  The Wigner amplitude can be negative or complex off diagonal and is
not subject to pointwise probability positivity
\cite{NeffFeldmeier2016}.

\subsection{Partial-wave decomposition}

With \(\Psi=\Psi_S+\Psi_D+\Psi_{\rm rel}\), the spectral amplitude decomposes
as
\begin{equation}
 \cS=\cS_{SS}+\cS_{SD}+\cS_{DS}+\cS_{DD}
 +\cS_{\rm rel}.
\label{eq:spectral-partial-wave-decomposition}
\end{equation}
The \(SD\) and \(DS\) terms are separately complex off forward and combine
under Hermiticity.  Tensor polarization, quadrupole structure, and many
spin--orbit observables receive signed interference contributions.  Replacing
\cref{eq:spectral-partial-wave-decomposition} by a scalar D-state probability
loses the phase and angular information needed for \(LT\), \(TT\), and
helicity-flip channels.

The controlled limits are:
\begin{align}
 \Psi_D\to0
 &\quad\Longrightarrow\quad
 \delta_T\rho=0
 \quad\text{for a nonrelativistic pure S-wave one-body density},
\label{eq:pure-s-tensor-zero}\\
 \pT\to0,\ y\to\tfrac12
 &\quad\Longrightarrow\quad
 R_M\to\bbone,
\label{eq:melosh-nonrel-limit}
\end{align}
subject to explicitly retained relativistic spin components.  A nonzero
pure-S tensor signal in a relativistic calculation must be traceable to those
components or to a many-body operator.

\subsection{Spectral normalization ledgers}

For an \(NN\) state normalized to \(Z_{\NN}\),
\begin{align}
 \sum_\lambda\int[\dd\Gamma_2]\,
 \rho_{\Lambda\Lambda;\lambda\lambda}^{p/D}
 &=Z_{\NN},
 &
 \sum_\lambda\int[\dd\Gamma_2]\,
 \rho_{\Lambda\Lambda;\lambda\lambda}^{n/D}
 &=Z_{\NN},
\label{eq:nucleon-number-spectral-ledger}\\
 \sum_{N=p,n}\sum_\lambda
 \int[\dd\Gamma_2]\,y_N
 \rho_{\Lambda\Lambda;\lambda\lambda}^{N/D}
 &=Z_{\NN}.
\label{eq:nucleon-momentum-spectral-ledger}
\end{align}
The first line counts one proton and one neutron in the \(NN\) sector.  The
second closes the sector plus momentum.  Full-state ledgers require the
corresponding \(NN\pi\), \(\Delta\Delta\), and compact-sector terms.

\section{Matched nuclear partonic operators}
\label{sec:nuclear-operators}

\subsection{Operator expansion in the nuclear theory}

The matched nuclear operator corresponding to a quark, antiquark, or gluon
GTMD has the expansion
\begin{equation}
 \widehat\cO_a^{D}
 =\widehat\cO_a^{(1)}
 +\widehat\cO_a^{(2)}
 +\widehat\cO_a^{(\pi)}
 +\widehat\cO_a^{(\Delta)}
 +\widehat\cO_a^{(\shortd)}
 +\widehat\cO_a^{(\coh)}
 +\cdots .
\label{eq:nuclear-operator-expansion}
\end{equation}
The superscripts identify operator topology, not independent probability
components:
\begin{itemize}
\item \(\widehat\cO^{(1)}\): matched nucleon-active operator, including
      induced one-body corrections in the selected representation;
\item \(\widehat\cO^{(2)}\): irreducible two-nucleon current or partonic
      operator;
\item \(\widehat\cO^{(\pi)}\): pion-active and pion-transition operators;
\item \(\widehat\cO^{(\Delta)}\): Delta-active and transition operators;
\item \(\widehat\cO^{(\shortd)}\): short-distance contact or compact-sector
      operator;
\item \(\widehat\cO^{(\coh)}\): coherent multiple-scattering operator in the
      factorization regime where such a reduction exists.
\end{itemize}
The same physical contribution may not appear in more than one term without
an explicit overlap subtraction.

\subsection{Complete operator identity}

A nuclear operator identity is
\begin{align}
 \mathfrak o_{a/D}=\bigl(&
 a,J,
 P_i,P_f,
 \gamma,R_{\rm color},
 \text{gluon link pair},\mathcal C_{f/d},
 \notag\\
 &\alpha,\beta,
 \text{active constituent},
 \mathfrak r_{\nuc},
 \mathfrak r_a,
 \cR_{\UV},\cR_{\rap},\cS_{\soft},
 \mu,\zeta,\mathfrak s,
 \text{rank},\mathfrak m_D
 \bigr),
\label{eq:nuclear-operator-identity}
\end{align}
where \(\alpha,\beta\) are source and target nuclear sectors.  The operator is
incomplete if a field required by the requested reduction is unspecified.
For example, a gluon operator with no ordered link pair or \(f/d\) color
status cannot enter a process map that requires those labels.

\subsection{One-body operator block}

In the \(NN\) sector the spectator-preserving one-body block is
schematically
\begin{equation}
 \widehat\cO_{a,\NN\to\NN}^{(1)}
 =\sum_{N=p,n}
 \int[\dd\Gamma_2]\,
 |N,s\rangle
 \widehat\cO_{a/N}
 \langle N,s|,
\label{eq:one-body-operator-block}
\end{equation}
with all helicity, recoil, and normalization indices retained.  The nucleon
operator in \cref{eq:one-body-operator-block} is the Volume~II/III parent, not
a named scalar TMD.

\subsection{Two-body and transition blocks}

A generic two-nucleon block has the form
\begin{equation}
 \widehat\cO_a^{(2)}
 =\sum_{A}
 C_{a,A}^{(2)}(x_D,\kT,\DT;\Lambda_D)
 \left[N^\dagger\Gamma_{A,1}N\right]
 \left[N^\dagger\Gamma_{A,2}N\right],
\label{eq:two-body-operator-basis}
\end{equation}
where \(A\) labels spin, isospin, derivative, color-singlet, Wilson, and
transverse-rank structure.  Its momentum kernel includes the transfer
partition in \cref{eq:two-body-transfer-partition}.  Transition blocks such
as \(\NN\leftrightarrow\NNpi\) or \(\NN\leftrightarrow\DD\) have explicit
source and target sectors and are accompanied by their Hermitian conjugates
when required.

\subsection{Matching conditions}

The coefficients in \cref{eq:nuclear-operator-expansion} are fixed by a
matching set
\begin{equation}
 \cM_{\rm match}
 =\left\{
 \langle X_f|\cO_a^{qg}|X_i\rangle
 =\langle X_f|\widehat\cO_a^{D}|X_i\rangle
 +\cO\!\left[(Q_{\nuc}/\Lambda_{\rm match})^{\nu+1}\right]
 \right\}
\label{eq:nuclear-operator-matching}
\end{equation}
for selected one- and two-hadron states, local moments, and symmetry
relations.  Matching may use lattice matrix elements, chiral constraints,
phenomenological form factors, or partonic observables only when their scheme
and scale are compatible.  A fit to one named deuteron TMD is not a matching
condition for the full operator basis.

\begin{requirement}[Shared coefficient propagation]
\label{req:shared-nuclear-coefficients}
A coefficient multiplying a two-body spin--isospin operator affects every
GTMD, GPD, PDF, current, and process reduction to which that operator
contributes.  It may not be duplicated as an independent coefficient for
each named TMD.
\end{requirement}

\subsection{Current consistency}

The electromagnetic, axial, energy--momentum, and partonic bilocal operators
must be constructed in the same unitary and regulator convention as
\(\cM_D^2\).  Chiral EFT calculations provide a concrete example: using
inconsistent regulators for forces and currents can violate the symmetry
relations the EFT is designed to preserve \cite{Krebs2020,Schiavilla2018}.
The same principle applies to the light-front nuclear theory.

\section{The full deuteron GTMD parent}
\label{sec:full-parent}

\subsection{Amplitude-level definition}

Let \(\alpha,\beta\) run over all retained nuclear sectors.  The full parent is
\begin{equation}
 W_{a/D,\Lambda'\Lambda}^{[J,\gamma]}
 (x_D,\kT,\DT)
 =\sum_{\alpha,\beta}
 \left\langle
 \Psi_{D,\Lambda'}^\alpha
 \middle|
 \widehat\cO_{a,\alpha\beta}^{[J,\gamma]}
 (x_D,\kT,\DT)
 \middle|
 \Psi_{D,\Lambda}^\beta
 \right\rangle.
\label{eq:full-deuteron-parent}
\end{equation}
For gluons, the ordered link pair and \(f/d\) color contraction are part of
\(\gamma\) and the operator identity.  The state member and nucleon
microscopic member remain correlated through the matching map.

\begin{principle}[Mechanisms are views of the parent]
\label{prin:mechanisms-views}
``Impulse,'' ``exchange current,'' ``pion,'' ``shadowing,'' ``short-range,''
and ``hidden color'' are labels on operator blocks, sector restrictions, or
controlled asymptotic reductions of \cref{eq:full-deuteron-parent}.  They are
not independently normalized functions to be summed after the fact.
\end{principle}

\subsection{Sector-diagonal and sector-changing contributions}

Equation~\eqref{eq:full-deuteron-parent} separates as
\begin{equation}
 W_{a/D}
 =\sum_\alpha W_{\alpha\alpha}
 +\sum_{\alpha<\beta}
 \left(W_{\alpha\beta}+W_{\beta\alpha}\right).
\label{eq:sector-decomposition-parent}
\end{equation}
A diagonal term \(W_{\NNpi,\NNpi}\) can contain an active pion or active
nucleon.  An off-diagonal term \(W_{\NN,\NNpi}\) requires an operator that
creates or absorbs the pion.  The existence of both sectors in the state is
not sufficient to authorize the transition matrix element.

\subsection{Hermiticity, parity, and link reversal}

For a Hermitian operator class,
\begin{equation}
 W_{a/D,\Lambda'\Lambda}^{[J,\gamma]}
 (x_D,\kT,\DT)
 =
 \left[
 W_{a/D,\Lambda\Lambda'}^{[J^\dagger,\gamma^{-1}]}
 (x_D,\kT,-\DT)
 \right]^*.
\label{eq:deuteron-parent-hermiticity}
\end{equation}
Parity acts through the declared light-front spin adapter.  If the nuclear
dynamics is invariant under the relevant antiunitary operation and the
process changes only the standard staple orientation, then
\begin{equation}
 W_{a/D,\odd}^{[\gamma_+]}
 =-W_{a/D,\odd}^{[\gamma_-]}.
\label{eq:deuteron-link-reversal}
\end{equation}
A failure of \cref{eq:deuteron-link-reversal} must be traced to an incomplete
state or current, an additional process link/color structure, a Glauber or
factorization obstruction, or a convention/numerical error.  It is not
repaired at the named-function level.

\subsection{Partonic and nuclear Wilson structures}

The partonic operator contains
\begin{equation}
 U_\gamma^{\rm parton}:
 \text{color transport associated with the hard factorization},
\label{eq:partonic-wilson-object}
\end{equation}
while coherent nuclear propagation is
\begin{equation}
 \cU_{\nuc}:
 \text{amplitude evolution among nuclear configurations}.
\label{eq:nuclear-propagation-object}
\end{equation}
If both descriptions include the same soft momentum region, the matched
parent contains a subtraction
\begin{equation}
 W_{a/D}^{\rm matched}
 =W_{a/D}^{\rm parton}
 +W_{a/D}^{\rm nuclear}
 -W_{a/D}^{\rm overlap}.
\label{eq:parton-nuclear-overlap-subtraction}
\end{equation}
The overlap is defined by a common asymptotic expansion, not by tuning the
sum to data.

\begin{requirement}[No nuclear phase multiplication]
\label{req:no-nuclear-phase-multiplication-v4}
A tensor-polarized T-odd deuteron function may not be generated by multiplying
an unpolarized nucleon TMD by a fitted nuclear phase or tensor ratio.  It must
be projected from the spin-resolved matrix element
\cref{eq:full-deuteron-parent} with the link-resolved nucleon operator and any
additional nuclear rescattering treated at amplitude level.
\end{requirement}

\subsection{Common-member covariance}

For one microscopic/nuclear member \(m\), all deuteron projections are
\begin{equation}
 F_A^{(m)}=\cP_A[W_{a/D}^{(m)}].
\label{eq:member-projection}
\end{equation}
The covariance is therefore generated by the same member ensemble,
\begin{equation}
 \operatorname{Cov}(F_A,F_B)
 =\left\langle
 (F_A^{(m)}-\bar F_A)(F_B^{(m)}-\bar F_B)
 \right\rangle_m,
\label{eq:common-member-covariance}
\end{equation}
retaining cross-flavor, proton/neutron, quark/gluon, TMD/GPD, and
nucleon/deuteron correlations.

\section{Impulse approximation as a controlled reduction}
\label{sec:impulse}

\subsection{Derivation of the one-body convolution}

Restrict \cref{eq:full-deuteron-parent} to the \(NN\) sector and the
spectator-preserving one-body operator.  Insert complete two-nucleon states,
use the recoil map in \cref{eq:active-intrinsic-recoil}, and define
\(z=x_D/y\), \(\ellT=\kT-z\pT\).  The result is
\begin{align}
 W_{a/D,\Lambda'\Lambda}^{[J,\gamma],\IA}
 (x_D,\kT,\DT)
 ={}&\sum_{N=p,n}
 \int_{x_D}^{1}\frac{\dd y}{y}
 \int\frac{\dd^2\pT}{(2\pi)^2}
 \sum_{\lambda',\lambda}
 \notag\\
 &\times
 \cS_{\Lambda'\Lambda;\lambda'\lambda}^{N/D}
 (y,\pT,\DT)
 W_{a/N,\lambda'\lambda}^{[J,\gamma]}
 \left(
 z,\ellT,\DNT=\DT
 \right)
 \cJ_{\IA}.
\label{eq:ia-gtmd-convolution}
\end{align}
The factor \(\cJ_{\IA}\) is fixed by the state and operator normalization;
it is not independently tuned in each polarization channel.  With a
different state normalization, the measure and \(\cJ_{\IA}\) transform
together.

Equation~\eqref{eq:ia-gtmd-convolution} retains:
\begin{itemize}
\item deuteron and active-nucleon helicity transitions;
\item proton and neutron identity;
\item quark flavor or gluon identity;
\item Wilson path and gluon color class;
\item the full external transfer and partonic transverse momentum;
\item nuclear and nucleon microscopic members.
\end{itemize}
Summing proton and neutron or projecting onto \(U,L,T,LL,LT,TT\) is a later
reduction.

\subsection{Forward TMD limit}

Setting \(\DT=0\) gives
\begin{align}
 \Phi_{a/D,\Lambda'\Lambda}^{[J,\gamma],\IA}
 (x_D,\kT)
 ={}&\sum_N
 \int_{x_D}^{1}\frac{\dd y}{y}
 \int\frac{\dd^2\pT}{(2\pi)^2}
 \sum_{\lambda',\lambda}
 \rho_{\Lambda'\Lambda;\lambda'\lambda}^{N/D}(y,\pT)
 \notag\\
 &\times
 \Phi_{a/N,\lambda'\lambda}^{[J,\gamma]}
 \left(\frac{x_D}{y},
 \kT-\frac{x_D}{y}\pT\right)
 \cJ_{\TMD}.
\label{eq:ia-tmd-convolution}
\end{align}
This is not a scalar smearing formula unless the nucleon and target spin
matrices are first restricted to a scalar channel.

\subsection{TMD-impact representation}

Using
\(
 \widetilde\Phi(x,\bM)=\int\dd^2\kT\,e^{i\bM\cdot\kT}\Phi(x,\kT)
\),
\cref{eq:ia-tmd-convolution} becomes
\begin{align}
 \widetilde\Phi_{a/D}^{\IA}(x_D,\bM)
 ={}&\sum_N
 \int_{x_D}^{1}\frac{\dd y}{y}
 \int\frac{\dd^2\pT}{(2\pi)^2}
 \rho^{N/D}(y,\pT)
 e^{i(x_D/y)\bM\cdot\pT}
 \notag\\
 &\times
 \widetilde\Phi_{a/N}
 \left(\frac{x_D}{y},\bM\right)
 \cJ_{\TMD},
\label{eq:ia-bspace-convolution}
\end{align}
with all spin indices implicit.  The nuclear phase in
\cref{eq:ia-bspace-convolution} is a kinematic boost phase; it is unrelated
to the partonic naive-T-odd Wilson phase.

\subsection{GPD, PDF, and current reductions}

Integrating \cref{eq:ia-gtmd-convolution} over \(\kT\) gives the zero-skewness
GPD convolution,
\begin{equation}
 H_{a/D}^{\IA}(x_D,0,t_D)
 =\int\dd^2\kT\,W_{a/D}^{\IA}(x_D,\kT,\DT),
\label{eq:ia-gpd-reduction}
\end{equation}
provided the nucleon staple operator is matched to the straight-link GPD
operator as required.  The forward collinear limit is
\begin{equation}
 f_{a/D}^{\IA}(x_D)
 =\int\dd^2\kT\,\Phi_{a/D}^{\IA}(x_D,\kT).
\label{eq:ia-pdf-reduction}
\end{equation}
Further \(x_D\) moments give local currents or energy--momentum matrix
elements only after the corresponding operator matching.

The required commuting square is
\begin{equation}
 \Red_D\circ\cN_{\IA}[W_{a/N}]
 =\cN_{\IA}^{\rm reduced}\circ\Red_N[W_{a/N}]
 +\delta_{\rm match}.
\label{eq:ia-reduction-square}
\end{equation}
The residual \(\delta_{\rm match}\) is reported separately as matching,
truncation, and numerical components.

\subsection{Support and number closure}

The convolution has \(0<x_D<1\), while the project coordinate
\(x_N=2x_D\) has support up to 2.  Fermi motion and compact sectors can
populate \(x_N>1\) without violating the canonical support in \(x_D\).

For valence quarks in the full deuteron,
\begin{equation}
 \int_0^1\dd x_D\,[u_D-\bar u_D]=3,
 \qquad
 \int_0^1\dd x_D\,[d_D-\bar d_D]=3,
\label{eq:deuteron-valence-ledger}
\end{equation}
and the partonic momentum sum is
\begin{equation}
 \sum_{a=q,\bar q,g}
 \int_0^1\dd x_D\,x_D f_{a/D}(x_D;\mu)=1.
\label{eq:deuteron-parton-momentum-sum}
\end{equation}
The impulse term alone closes these relations only within the \(NN\) sector
and its matched one-body operator.  Pion, Delta, compact, and many-body
operator contributions must be included in the full ledger.

\subsection{Completely positive impulse reduction}

At fixed nuclear kinematics, define an amplitude embedding
\begin{equation}
 V_D:\Hil_N\longrightarrow\Hil_D\otimes\Hil_{\spec}
\label{eq:nuclear-amplitude-embedding}
\end{equation}
from the deuteron wave function.  If the spectator is unobserved and no
coherence with another sector remains, the reduced map is
\begin{equation}
 \cE_D(\rho_N)
 =\Tr_{\spec}\left[V_D\rho_NV_D^\dagger\right]
 =\sum_\alpha K_\alpha\rho_NK_\alpha^\dagger.
\label{eq:impulse-cp-map}
\end{equation}
This proves complete positivity of the reduced impulse channel.  It does not
justify representing coherent shadowing, exchange currents, or
\(NN\leftrightarrow NN\pi\) interference by independent Kraus maps.

\section{Spin recoupling and the complete spin-1 projection system}
\label{sec:spin-recoupling}

\subsection{Why the nucleon helicity matrix is indispensable}

The deuteron target projector and active-parton projector act only after the
nuclear convolution.  For a named channel \(A\),
\begin{equation}
 F_{A,a/D}
 =\Tr_{\Lambda,\lambda_a}
 \left[
 P_{A}^{(D,a)}
 W_{a/D}
 \right].
\label{eq:deuteron-tmd-projector}
\end{equation}
Replacing \(W_{a/N,\lambda'\lambda}\) in
\cref{eq:ia-gtmd-convolution} by a spin-averaged scalar before convolution
removes terms in which vector-polarized nucleon distributions feed tensor
polarization through deuteron spin recoupling.

\begin{requirement}[Late projection]
\label{req:late-spin-projection}
The active-nucleon helicity indices and deuteron helicity indices remain
explicit through nuclear composition.  Projection onto \(U,L,T,LL,LT,TT\)
and named parton polarization occurs only after all operator blocks needed by
the requested observable are assembled.
\end{requirement}

\subsection{Irreducible nuclear response}

Apply \cref{eq:spin1-irrep-coefficients} to the deuteron helicity matrix:
\begin{equation}
 W_{a/D;M}^{(K)}
 =\sum_{\Lambda',\Lambda}
 (T_M^{(K)})_{\Lambda\Lambda'}^*
 W_{a/D,\Lambda'\Lambda}.
\label{eq:deuteron-irrep-projection}
\end{equation}
The physical mapping is
\begin{equation}
 K=0\leftrightarrow U,
 \qquad
 K=1\leftrightarrow L,T,
 \qquad
 K=2\leftrightarrow LL,LT,TT.
\label{eq:irrep-physical-map}
\end{equation}
The \(M\) dependence fixes the transverse \(SO(2)\) weight that must combine
with the partonic tensor rank and orbital interference.

\subsection{S--D and orbital selection}

A nuclear contribution is nonzero only when total light-front angular
momentum closes:
\begin{equation}
 \Delta\Lambda_D
 =\Delta\lambda_N
 +\Delta L_{z,\nuc}
 +m_{\cO,\nuc}.
\label{eq:nuclear-jz-selection}
\end{equation}
The active nucleon operator then has its own partonic balance inherited from
Volume~II.  The complete selection rule is the composition of the nuclear and
partonic intertwiners, not one plot-level rank factor.

In the lowest \(NN\) component:
\begin{itemize}
\item \(SS\) supplies the dominant unpolarized scalar response;
\item \(SD+DS\) supplies signed quadrupole and many tensor interferences;
\item \(DD\) supplies additional tensor and high-momentum strength;
\item Melosh and additional relativistic components permit spin couplings
      absent in the nonrelativistic limit.
\end{itemize}
Disabling either amplitude in a required interference block must remove the
corresponding harmonic.

\subsection{Tensor PDF and the \texorpdfstring{\(b_1\)}{b1} anchor}

For any partonic species,
\begin{equation}
 \delta_T f_{a/D}(x_D)
 =f_{a/D}^{\Lambda=0}(x_D)
 -\frac12\left[
 f_{a/D}^{\Lambda=+1}(x_D)
 +f_{a/D}^{\Lambda=-1}(x_D)
 \right].
\label{eq:tensor-pdf-definition}
\end{equation}
At leading twist and leading order,
\begin{equation}
 b_1(x_D,Q^2)
 =\frac12\sum_q e_q^2
 \left[
 \delta_T q_D(x_D,Q^2)
 +\delta_T\bar q_D(x_D,Q^2)
 \right]
 +\cO(\alpha_s,M_D^2/Q^2).
\label{eq:b1-relation}
\end{equation}
With the project convention \(S_{LL}^{\Lambda=0}=-1\),
\begin{equation}
 f_{1LL}=-\frac23\,\delta_T f_1.
\label{eq:f1ll-tensor-adapter}
\end{equation}
The adapter is applied only after the convention-independent helicity
difference is computed.  The impulse term from \(S\)-\(D\) structure is a
baseline, not a guarantee of agreement with data
\cite{HoodbhoyJaffeManohar1989,HERMES2005,CosynDongKumanoSargsian2017}.

\subsection{Gluon double-helicity-flip structure}

A spin-1 target supports gluon helicity-flip-two matrix elements with no
leading one-body spin-\(\tfrac12\) nucleon analogue.  In a deuteron they can
arise from nuclear orbital transfer, coherent two-nucleon operators, or
compact/non-nucleonic sectors.  The model must therefore preserve the
operator-level distinction
\begin{equation}
 W_{g/D}^{\Delta\lambda_g=2}
 \neq
 C\,W_{g/D}^{\rm unpolarized}.
\label{eq:gluon-transversity-not-rescaling}
\end{equation}
A nonzero result identifies specific angular-momentum and sector blocks; a
zero impulse baseline is not an exact all-mechanism zero.

\subsection{Forward positivity domain}

If the complete forward correlator is represented as a sum over cut
amplitudes,
\begin{equation}
 \Phi_{AB}(x_D,\kT)
 =\sum_X \cA_{AX}\cA_{BX}^*,
\label{eq:deuteron-forward-gram}
\end{equation}
then it is positive semidefinite in the joint parton--deuteron helicity
space.  Positivity is imposed on the total physical correlator at the layer
where \cref{eq:deuteron-forward-gram} is valid.  Individual shadowing,
exchange, or subtraction terms need not be positive, and off-forward GTMDs
and Wigner distributions are not pointwise probability matrices.

\section{Consistent currents, form factors, and local moments}
\label{sec:currents-formfactors}

\subsection{Spin-1 electromagnetic current}

For on-shell deuterons, current conservation and discrete symmetries permit
three electromagnetic form factors.  With polarization vectors
\(\epsilon^\alpha(P_i,\Lambda)\) and
\(\epsilon^{*\beta}(P_f,\Lambda')\), one convention is
\begin{align}
 \langle P_f,\Lambda'|J^\mu|P_i,\Lambda\rangle
 ={}&-\epsilon_f^{*\alpha}\epsilon_i^\beta
 \Bigg\{
 (P_f+P_i)^\mu
 \left[
 g_{\alpha\beta}G_1(Q^2)
 -\frac{\Delta_\alpha\Delta_\beta}{2M_D^2}G_3(Q^2)
 \right]
 \notag\\
 &\hspace{28mm}
 +\left(g^\mu{}_{\alpha}\Delta_\beta
 -g^\mu{}_{\beta}\Delta_\alpha\right)G_2(Q^2)
 \Bigg\}.
\label{eq:spin1-current-decomposition}
\end{align}
Defining \(\eta=Q^2/(4M_D^2)\),
\begin{align}
 G_C&=G_1+\frac23\eta G_Q,
 &
 G_M&=G_2,
 &
 G_Q&=G_1-G_2+(1+\eta)G_3.
\label{eq:physical-deuteron-formfactors}
\end{align}
The static conditions are
\begin{equation}
 G_C(0)=1,
 \qquad
 G_M(0)=\frac{M_D}{M_N}\mu_D,
 \qquad
 G_Q(0)=M_D^2Q_D
\label{eq:static-formfactor-normalization}
\end{equation}
in the declared magnetic-moment and quadrupole conventions.

\subsection{One- and two-body current pair}

The nuclear current is
\begin{equation}
 J_D^\mu
 =J_{(1)}^\mu+J_{(2)}^\mu+J_{(\pi)}^\mu
 +J_{(\Delta)}^\mu+J_{(\shortd)}^\mu+\cdots.
\label{eq:nuclear-current-expansion}
\end{equation}
Current conservation requires
\begin{equation}
 \Delta_\mu\langle D_f|J_D^\mu|D_i\rangle=0
\label{eq:deuteron-current-conservation}
\end{equation}
within the declared truncation.  A one-body current combined with a
Hamiltonian containing exchange or eliminated-sector interactions generally
does not satisfy the same Ward identity without induced current terms.

\begin{requirement}[Current-family consistency]
\label{req:current-family-consistency}
Every nuclear member identifies the Hamiltonian, regulator, one-body current,
two-body current, partonic operator, and matching order as one family.  A
current imported from another interaction is an explicit hybrid sensitivity,
not the central matched prediction.
\end{requirement}

\subsection{Angular-condition closure}

The residual in \cref{eq:angular-condition} is evaluated for each nuclear
member and current order.  Different prescriptions for extracting
\(G_C,G_M,G_Q\) from an overcomplete set of light-front helicity amplitudes
must agree as \(\Delta_{\rm ang}\to0\).  The spread among prescriptions at
finite truncation is reported as a rotational/current defect and is not
folded into an unrelated statistical band.

\subsection{GTMD moments and local operators}

The matched straight-link GPD reduction of the deuteron GTMD satisfies
schematically
\begin{align}
 \int\dd x_D\,H_{q/D}(x_D,0,t)
 &\longrightarrow
 \langle D_f|J_q^+|D_i\rangle,
\label{eq:vector-current-moment}\\
 \sum_{a=q,\bar q,g}
 \int\dd x_D\,x_D H_{a/D}(x_D,0,t)
 &\longrightarrow
 \langle D_f|T_a^{++}|D_i\rangle.
\label{eq:emt-moment}
\end{align}
The local-current route and the GTMD-moment route must use the same matched
operator basis.  Agreement is a nontrivial closure test of the nuclear
recoil, current, and sector structure.

\subsection{Static and elastic validation set}

Before using the state for partonic predictions, the nuclear theory must
reproduce within declared uncertainties:
\begin{equation}
 B_D,
 \quad
 \mu_D,
 \quad
 Q_D,
 \quad
 G_C(Q^2),G_M(Q^2),G_Q(Q^2),
 \quad
 \Delta_{\rm ang}(Q^2).
\label{eq:elastic-validation-set}
\end{equation}
The long-distance \(S/D\) phenomenology of high-quality interactions is an
important benchmark, but success in elastic observables does not by itself
validate partonic two-body operators.

\section{Coherent small-x amplitudes and shadowing}
\label{sec:coherent-shadowing}

\subsection{Amplitude-level origin}

At small \(x_D\), the coherence length of the hard fluctuation can exceed the
nucleon separation and permit scattering from both nucleons.  The deuteron
matrix element then contains a coherent two-step amplitude rather than an
incoherent modification of one nucleon distribution.  Schematically,
\begin{align}
 W_{a/D}^{(2\,\coh)}
 \sim{}&
 \sum_{X}
 \int[\dd\Gamma_2][\dd\Gamma_2']
 \Psi_D^*(\Gamma_2')
 \cA_{N_2\to X}^{a,\rm diff}
 \notag\\
 &\times
 e^{i\varphi_{\coh}(x_D;\Gamma_2',\Gamma_2)}
 \cA_{X\to N_1}^{a,\rm diff}
 \Psi_D(\Gamma_2),
\label{eq:coherent-double-scattering}
\end{align}
where \(X\) is the diffractively produced intermediate state and
\(\varphi_{\coh}\) carries the longitudinal and transverse propagation
phase.  Leading-twist nuclear shadowing connects this amplitude to nucleon
diffractive parton distributions in the kinematic domain where the
factorization assumptions hold \cite{FrankfurtGuzeyStrikman2012}.

\subsection{Helicity-resolved diffractive kernel}

The elementary diffractive input is a matrix
\begin{equation}
 \cA_{a;\lambda_N'\lambda_N}^{\rm diff}
 (x_{\mathbb P},\beta,t;J,\gamma),
\label{eq:helicity-diffractive-amplitude}
\end{equation}
not one unpolarized ratio.  The coherent nuclear kernel acts on the full
nucleon and deuteron helicity structure:
\begin{equation}
 \cK_{a;\lambda_1'\lambda_2';\lambda_1\lambda_2}^{\coh}
 =\sum_X
 \cA_{a;\lambda_1'\lambda_1}^{\rm diff}
 \cP_X
 \cA_{a;\lambda_2'\lambda_2}^{\rm diff}
 \times\text{propagation phase}.
\label{eq:coherent-helicity-kernel}
\end{equation}
Vector and tensor shadowing are projections of
\cref{eq:coherent-helicity-kernel}.  They cannot be obtained by multiplying
the unpolarized shadowing correction by a target-polarization coefficient.

\begin{requirement}[Independent polarized diffraction]
\label{req:independent-polarized-diffraction}
An unpolarized diffractive PDF or amplitude constrains only the unpolarized
part of the coherent kernel.  Polarized and tensor components require their
own amplitudes, symmetry relations, data, lattice input, or explicitly named
model priors.  Their uncertainty is not the unpolarized error band copied
between channels.
\end{requirement}

\subsection{Shadowing and antishadowing}

The coherent correction is negative in many inclusive unpolarized channels,
but neither sign nor magnitude is universal for all helicity and color
operators.  Antishadowing is represented by a distinct matched operator or
unitarity/moment-restoring mechanism,
\begin{equation}
 W_{a/D}^{\coh,\rm total}
 =W_{a/D}^{\rm shadow}
 +W_{a/D}^{\rm antishadow},
\label{eq:shadow-antishadow}
\end{equation}
with separate provenance.  It is not an arbitrary positive function added to
force the momentum sum rule; the sum rule is a validation constraint on the
combined state and operator.

\subsection{Coherence and the partonic soft sector}

A coherent nuclear exchange can overlap the Glauber or soft region already
represented by the Volume~III Wilson and soft factors.  The overlap term in
\cref{eq:parton-nuclear-overlap-subtraction} is therefore mandatory whenever
the asymptotic expansions coincide.  The subtraction record contains:
\begin{equation}
\begin{aligned}
 \mathfrak s_{\rm overlap}=\bigl(&
 \text{momentum region},
 \text{Wilson geometry},\\
 &\text{color representation},
 \text{nuclear channel},\\
 &\text{expansion order},
 \text{counterterm}
 \bigr).
\end{aligned}
\label{eq:overlap-subtraction-record}
\end{equation}

\subsection{Multiple-scattering convergence}

The coherent series
\begin{equation}
 W^{\coh}=W^{(2)}+W^{(3)}+\cdots
\label{eq:coherent-series}
\end{equation}
is organized by the number of resolved scatterings or by an equivalent
unitarized representation.  For the deuteron, the leading nontrivial
nucleonic term is double scattering, while higher terms can arise through
explicit additional constituents or internal rescattering structure.  The
truncation order and unitarization prescription are separate member fields.

\section{Binding, Fermi motion, off-shell representations, and short range}
\label{sec:binding-offshell-src}

\subsection{Joint spectral--partonic amplitude}

Binding and Fermi motion enter through the exact nuclear state and recoil
map.  A covariant spectral representation may be written schematically as
\begin{equation}
 W_{a/D}
 =\int\dd^4p\,
 \cS_D(p+\tfrac\Delta2,p-\tfrac\Delta2)
 W_{a/N}(p,\Delta)
 +W_{a/D}^{(2+)}.
\label{eq:covariant-spectral-representation}
\end{equation}
The separation between an off-shell continuation of \(W_{a/N}\) and the
many-body term \(W_{a/D}^{(2+)}\) depends on the representation, as shown in
\cref{prop:no-offshell-observable}.  Only their matched sum is physical.

\subsection{Controlled off-shell sensitivity}

For diagnostic purposes one may expand
\begin{equation}
 W_{a/N}(p^2)
 =W_{a/N}(M_N^2)
 +(p^2-M_N^2)
 \left.\frac{\partial W_{a/N}}{\partial p^2}\right|_{M_N^2}
 +\cdots,
\label{eq:offshell-expansion}
\end{equation}
but the derivative and compensating two-body operator belong to one scheme.
The project may report the size of the chosen representation's one-body
term, provided it is not called an observable and the complete result is
shown to be invariant under an allowed transformation.

\subsection{Short-range correlations and large x}

High relative momentum in the \(NN\) state and explicit short-distance
sectors can populate the large-\(x_N\) region.  The nuclear theory must
distinguish:
\begin{enumerate}
\item high-momentum components generated by the chosen \(NN\) interaction;
\item induced two-body partonic operators at the same resolution;
\item explicit cluster or compact-sector amplitudes;
\item perturbative high-\(k_T\) tails supplied later by matching and the
      observable \(Y\) term.
\end{enumerate}
Counting all four as independent additive ``SRC enhancements'' is forbidden.

\subsection{EMC-like behavior}

An EMC-like modification at intermediate \(x_D\) may emerge from binding,
Fermi motion, changed operator coefficients, explicit mesonic or compact
sectors, or short-distance two-body operators.  The model should not begin
with a universal multiplicative EMC ratio for every TMD.  It should compare
candidate mechanisms at the operator and observable levels and reserve at
least one flavor/spin family for discrimination.

\subsection{Tagged observables as spectral diagnostics}

Spectator tagging resolves \(y\), \(\pT\), and spectator spin information and
therefore tests the same spectral amplitude used in
\cref{eq:ia-gtmd-convolution}.  A model that fits inclusive nuclear ratios but
fails tagged momentum and polarization dependences has not validated its
nuclear state.  The tagged formalism is developed further in
\cref{sec:tagged-observables}.

\section{The spin-resolved \texorpdfstring{\(NN\pi\)}{NNpi} sector}
\label{sec:nnpi}

\subsection{Three-body light-front amplitude}

The \(NN\pi\) component is
\begin{align}
 |\Psi_{\NNpi}^{\Lambda}\rangle
 ={}&\sum_{\lambda_1,\lambda_2,a_\pi}
 \int[\dd\Gamma_3]\,
 \Psi_{\Lambda;\lambda_1\lambda_2,a_\pi}^{\NNpi}
 (y_1,y_2,y_\pi;\bm p_{1T},\bm p_{2T},\bm p_{\pi T})
 \notag\\
 &\times
 |N_1N_2\pi^{a_\pi}\rangle,
\label{eq:nnpi-state}
\end{align}
with
\begin{equation}
 y_1+y_2+y_\pi=1,
 \qquad
 \bm p_{1T}+\bm p_{2T}+\bm p_{\pi T}=0.
\label{eq:nnpi-closure}
\end{equation}
Its norm is \(Z_{\NNpi}\), and its spin, isospin, orbital, regulator, and
matching origin are explicit.  A scalar pion momentum distribution is only a
marginal of \cref{eq:nnpi-state}.

\subsection{Operator content}

The \(NN\pi\) sector contributes through:
\begin{align}
 W_{\NNpi,\NNpi}
 ={}&W_{\rm active\ N}
 +W_{\rm active\ \pi}
 +W_{\rm two\ body}
 +W_{\rm interference},
\label{eq:nnpi-diagonal-content}\\
 W_{\NN,\NNpi}+W_{\NNpi,\NN}
 ={}&\text{pion-emission/absorption transition operators}.
\label{eq:nn-nnpi-transition}
\end{align}
Each term has a matched operator identity.  The active-pion term consumes a
pion quark/gluon correlator in the same QCD scheme; the transition terms are
not generated by multiplying the pion PDF by a deuteron splitting function.

\subsection{Sullivan convolution as a reduction}

When the three-body amplitude factorizes in the Sullivan regime and the
observable does not resolve transition coherence, the pion-active diagonal
term reduces to
\begin{equation}
 F_{a/D}^{(\pi)}(x_D,\kT)
 =\int_{x_D}^{1}\frac{\dd y_\pi}{y_\pi}
 \int\dd^2\qT\,
 f_{\pi/D}(y_\pi,\qT;\Lambda',\Lambda)
 F_{a/\pi}
 \left(\frac{x_D}{y_\pi},
 \kT-\frac{x_D}{y_\pi}\qT\right)
 +\delta F_{\rm trans}.
\label{eq:sullivan-reduction}
\end{equation}
The spin-resolved splitting amplitude determines the deuteron tensor
structure.  The scalar function \(f_{\pi/D}\) is insufficient for general
\(LT\), \(TT\), and T-odd channels.

\subsection{Internal and exchange pion subtraction}

The microscopic nucleon state may already contain \(|N\pi\rangle\) or more
general chiral Fock sectors.  The nuclear theory also contains pions exchanged
between nucleons.  Let \(\mathsf T_{\pi}^{\rm int}\) and
\(\mathsf T_{\pi}^{\rm exch}\) denote asymptotic expansion maps that extract,
respectively, the internal-nucleon and inter-nucleon low-momentum regions.
They are not orthogonal projectors and do not define probabilities.  In their
shared domain the two descriptions must possess one common expansion,
\begin{equation}
 \mathsf T_{\pi}^{\rm exch}W^{\pi,\rm nucleon}
 =\mathsf T_{\pi}^{\rm int}W^{\pi,\rm nuclear}
 \equiv W^{\pi,\rm overlap}
 +\mathcal O(\delta_{\pi,\rm match}),
\label{eq:pion-region-overlap}
\end{equation}
where \(\delta_{\pi,\rm match}\) is the declared matching remainder.  The
matched result is then
\begin{equation}
 W^{\pi,\rm matched}
 =W^{\pi,\rm nucleon}
 +W^{\pi,\rm nuclear}
 -W^{\pi,\rm overlap}.
\label{eq:pion-double-counting-subtraction}
\end{equation}
The overlap term is fixed by the common asymptotic expansion, not by a fitted
normalization.  The pion contribution is therefore included once.

\subsection{Light-front Hamiltonian EFT status}

Light-front Hamiltonian EFT with explicit pion Fock sectors provides a
promising route to nonperturbative longitudinal momentum distributions and
has been extended conceptually toward the deuteron
\cite{CaoEtAl2026}.  For the present program it becomes a central replacement
only when it supplies the spin-resolved three-body deuteron amplitude,
operator matching, and convergence required above.  A longitudinal pion
marginal alone remains comparison information.

\section{Delta, six-quark, hidden-color, and other compact sectors}
\label{sec:compact-sectors}

\subsection{Delta--Delta sector}

The allowed isoscalar \(\Delta\Delta\) component is
\begin{equation}
 |\Psi_{\DD}^{\Lambda}\rangle
 =\sum_{\lambda_1,\lambda_2,I_1,I_2}
 \int[\dd\Gamma_{\DD}]\,
 \Psi_{\DD}^{\Lambda}
 |\Delta_1\Delta_2\rangle_{I=0,J=1}.
\label{eq:deltadelta-state}
\end{equation}
Its partonic prediction requires Delta quark/gluon operators and
\(NN\leftrightarrow\Delta\Delta\) transition operators in the same matching
scheme.  A quoted \(\Delta\Delta\) probability without these amplitudes does
not define a GTMD contribution.

\subsection{Six-quark color space}

Six quarks can couple to an overall color singlet through more than the
proton--neutron singlet--singlet channel.  Let
\begin{equation}
 \Hil_{\sixq}^{\rm singlet}
 =\Hil_{1\otimes1}
 \oplus\Hil_{\rm hidden\ color}
\label{eq:sixq-color-decomposition}
\end{equation}
be an orthogonal color-coupling basis at the chosen resolution.  Hidden color
is a basis statement within the compact six-quark sector; its probability is
basis and resolution dependent.  Physical predictions require the full
six-quark amplitude and operator matrix elements.

\subsection{Compact-sector operator requirements}

A nonzero central compact-sector contribution must specify:
\begin{enumerate}
\item normalized six-quark or cluster amplitudes;
\item explicit \(u,d,\bar q,g\) content;
\item quark and gluon color couplers;
\item parton and target helicities;
\item orbital components sufficient for the claimed transverse ranks;
\item nonzero-transfer recoil and support;
\item Wilson paths and gluon \(f/d\) structures;
\item transition operators to the \(NN\) and other sectors;
\item regulator and matching convergence.
\end{enumerate}
Absent these items, the sector remains a zero-centered sensitivity or an
external benchmark.  The existing hidden-color descriptions of \(b_1\) are
valuable phenomenological tests, but do not by themselves supply the full
spin-resolved GTMD parent \cite{Miller2014,MillerHiddenColor2014}.

\subsection{Exotic and null-test channels}

Gluon double-helicity-flip and selected tensor observables are especially
sensitive to coherent or compact structure because their lowest one-body
nucleon contribution can vanish or be suppressed.  Such observables are
powerful diagnostics only if the model keeps the exact null hypothesis:
\begin{equation}
 \text{one-body nucleon baseline}=0
 \quad\not\Rightarrow\quad
 \text{full deuteron observable}=0.
\label{eq:null-not-full-zero}
\end{equation}
The first nonzero operator block must be identified explicitly.

\subsection{Short-distance contacts versus explicit compact sectors}

A compact six-quark sector and a short-distance two-nucleon contact operator
can represent the same low-resolution physics.  The provenance graph must
therefore contain a replacement or matching two-cell,
\begin{equation}
 \sixq\ \text{explicit}
 \quad\longleftrightarrow\quad
 \widehat\cO^{(\shortd)}\ \text{with compact sector eliminated},
\label{eq:compact-contact-equivalence}
\end{equation}
and reject their simultaneous unrestricted use.

\section{Coherent amplitudes and completely positive reductions}
\label{sec:cp-reductions}

\subsection{Amplitude embedding is primary}

Let the complete nuclear amplitude define an isometry or contraction
\begin{equation}
 V:\Hil_{\rm resolved}\longrightarrow
 \Hil_{\rm observed}\otimes\Hil_{\rm unobserved}.
\label{eq:stinespring-embedding}
\end{equation}
When the observable is insensitive to the unobserved subsystem, the reduced
map
\begin{equation}
 \cE(\rho)
 =\Tr_{\rm unobserved}\left[V\rho V^\dagger\right]
\label{eq:stinespring-reduction}
\end{equation}
is completely positive.  A Kraus representation follows by choosing an
orthonormal basis of the unobserved space.

\begin{proposition}[Derived complete positivity]
\label{prop:derived-cp}
If \(V\) is the amplitude map for a complete set of retained unresolved
outcomes, then \cref{eq:stinespring-reduction} is completely positive and is
trace preserving when \(V^\dagger V=\bbone\).  If only selected outcomes are
retained, it is trace nonincreasing.
\end{proposition}

\begin{proof}
For any auxiliary space, \((\cE\otimes\Id)(X)\) is the partial trace of
\((V\otimes\bbone)X(V^\dagger\otimes\bbone)\), which is positive for positive
\(X\).  The trace statement follows from cyclicity and \(V^\dagger V\).
\end{proof}

\subsection{Where a CP map is valid}

Examples include:
\begin{itemize}
\item the forward impulse density after tracing an unobserved spectator;
\item detector or acceptance coarse graining after amplitudes are formed;
\item an inclusive sum over mutually orthogonal final nuclear channels;
\item selected stochastic response approximations derived from an amplitude
      embedding.
\end{itemize}

\subsection{Where it is not valid}

A CP map may not replace:
\begin{itemize}
\item coherent shadowing before the diffractive amplitudes are added;
\item \(NN\leftrightarrow NN\pi\) interference;
\item exchange-current interference with the impulse amplitude;
\item off-forward sector coherence;
\item Wilson-link phases whose sign follows from antiunitary reversal;
\item subtraction terms, which can be negative and are not physical channels.
\end{itemize}
Forcing these terms into independent positive maps can erase the interference
that produces tensor observables.

\subsection{Relation to the accepted boundary}

Ordered completely positive response maps in the accepted phenomenological
boundary remain useful reduced approximations and regression objects.  In the
microscopic hierarchy they must be derived from, or calibrated against, an
explicit amplitude-level calculation.  They cannot be treated as proof that
the underlying coherent nuclear dynamics has been solved.

\section{Tagged spectator observables as differential validation}
\label{sec:tagged-observables}

\subsection{Tagged process and variables}

Consider
\begin{equation}
 e(l)+D(P_D,\rho_D)
 \longrightarrow
 e(l')+N_s(p_s,\lambda_s)+X,
\label{eq:tagged-process}
\end{equation}
where the detected nucleon is in the target-fragmentation region.  Define
\begin{equation}
 y_s=\frac{p_s^+}{P_D^+},
 \qquad
 y=1-y_s,
 \qquad
 \pT=-\bm p_{sT}
\label{eq:tagged-variables}
\end{equation}
in the deuteron transverse rest frame.  The tagged momentum fixes the nuclear
configuration rather than integrating it into an inclusive smearing function.

\subsection{Impulse tagged tensor}

In impulse approximation the tagged hadronic tensor is
\begin{align}
 W_{D,\tagg}^{\mu\nu}
 (P_D,q;p_s)_{\Lambda'\Lambda}
 ={}&\sum_{N=p,n}
 \sum_{\lambda',\lambda}
 \cS_{\Lambda'\Lambda;\lambda'\lambda}^{N/D}
 (y,\pT,0)
 \notag\\
 &\times
 W_N^{\mu\nu}(p_N,q)_{\lambda'\lambda}
 \cJ_{\tagg}
 +W_{D,\tagg}^{\mu\nu,\rm FSI}
 +W_{D,\tagg}^{\mu\nu,\rm nonIA}.
\label{eq:tagged-hadronic-tensor}
\end{align}
The same spectral amplitude appears in
\cref{eq:ia-gtmd-convolution}; tagged data therefore test the state and spin
recoupling used in the inclusive and TMD calculations.

General spin-1 tagged observables require vector and tensor target density
matrices, including transverse components.  Recent light-front formulations
give a complete covariant cross-section and spectator-tagging decomposition
for arbitrary spin-1 polarization
\cite{CosynWeiss2020,CosynWeiss2026I,CosynWeiss2026II}.

\subsection{Inclusive recovery}

Summing spectator spin and integrating over the complete spectator phase
space must recover the corresponding inclusive impulse tensor:
\begin{equation}
 \sum_{\lambda_s}
 \int[\dd\Gamma_s]\,
 W_{D,\tagg}^{\mu\nu,\IA}
 =W_D^{\mu\nu,\IA}.
\label{eq:tagged-inclusive-sumrule}
\end{equation}
The equality tests the flux factor, wave-function normalization, active-slot
multiplicity, and spin-density algebra.  Detector cuts are applied only after
this theoretical closure has been established.

\subsection{On-shell pole and neutron extraction}

For proton tagging, the active neutron virtuality is
\begin{equation}
 t_N=(P_D-p_s)^2,
\label{eq:active-neutron-virtuality}
\end{equation}
and the impulse amplitude contains the nucleon pole at \(t_N=M_N^2\).  Near
the pole,
\begin{equation}
 W_{D,\tagg}^{\IA}
 =\frac{R_D(p_s)}{(t_N-M_N^2)^2}
 W_N^{\rm on}
 +\text{less singular terms}.
\label{eq:tagged-pole-form}
\end{equation}
Analytic extrapolation of
\((t_N-M_N^2)^2W_{D,\tagg}\) to the pole isolates the on-shell nucleon residue
under the assumptions of the method.  Final-state interactions and many-body
terms are less singular and can be constrained by the spectator-momentum
dependence \cite{CosynSargsianWeiss2016,CosynWeiss2020}.

\begin{requirement}[Pole extrapolation is a validation, not a fit shortcut]
\label{req:pole-validation}
The model must predict the pole position, residue normalization, spin
structure, and approach to the pole from its own spectral amplitude.  It may
not fit an arbitrary spectator factor separately in each tagged observable.
\end{requirement}

\subsection{Tensor tagging and D-wave selection}

Tensor-polarized tagged observables can isolate regions with enhanced
D-wave and spin--orbit content.  Their dependence on spectator momentum and
azimuth tests the \(SD,DS,DD\) blocks of
\cref{eq:spectral-partial-wave-decomposition} directly.  A successful
inclusive \(b_1\) curve with incorrect tagged tensor asymmetries signals a
misidentified nuclear mechanism or spin recoupling.

\subsection{Tagged GTMD extension}

A future off-forward tagged parent retains both the spectator momentum and
the external transfer,
\begin{equation}
 W_{a/D,\tagg}
 (x_D,\kT,\DT;p_s,\lambda_s),
\label{eq:tagged-gtmd-parent}
\end{equation}
and can separate active-nucleon and coherent transfer patterns.  The first
implementation may remain forward in the tagged channel, but its types must
not collapse \(p_{sT}\), \(\pT\), \(\DT\), or \(\DNT\).

\section{Global ledgers, sum rules, and no-double-counting closure}
\label{sec:sumrules}

\subsection{State and parton ledgers}

For every nuclear member, the following quantities are evaluated from the
same state and matched operators:
\begin{enumerate}
\item state normalization and sector probabilities;
\item baryon number and electric charge;
\item hadronic plus-momentum sharing among nucleons, pions, Deltas, and
      compact sectors;
\item partonic momentum sharing among quarks, antiquarks, and gluons;
\item vector and tensor spin moments;
\item local electromagnetic and energy--momentum moments;
\item tagged-to-inclusive closure.
\end{enumerate}
The two momentum ledgers are different: the hadronic ledger partitions the
nuclear EFT degrees of freedom, whereas the partonic ledger resolves the
matched quark/gluon content inside all sectors.  They must be connected by
matching, not added as separate normalizations.

\subsection{Tensor number and conditional sum rules}

The helicity-state normalization implies
\begin{equation}
 \int\dd x_D\,
 \delta_T\left[q_D(x_D)-\bar q_D(x_D)\right]=0
\label{eq:tensor-valence-number}
\end{equation}
for each conserved flavor charge when all operator and endpoint terms are
included.  Additional \(b_1\) or tensor-PDF sum rules require assumptions
about tensor-polarized sea distributions, small-\(x\) convergence, and
higher-twist terms.  They are implemented with their hypotheses attached,
not as unconditional hard constraints
\cite{Kumano2014,KumanoSong2021}.

\subsection{Momentum and antishadowing closure}

The full partonic momentum sum in \cref{eq:deuteron-parton-momentum-sum}
contains impulse, meson, coherent, transition, and compact contributions.
The closure residual is
\begin{equation}
 \eps_{P^+}
 =1-
 \sum_{a=q,\bar q,g}
 \int_0^1\dd x_D\,x_D f_{a/D}(x_D).
\label{eq:momentum-closure-residual}
\end{equation}
A nonzero residual diagnoses missing matching or support.  It is not removed
by renormalizing the final distributions after composition.

\subsection{No-double-counting graph}

The nuclear provenance graph contains at least the relations
\begin{equation}
\begin{gathered}
 \{\mathrm{DERIVES\_FROM},\ \mathrm{PROJECTS\_TO},\
   \mathrm{ADDS\_TO},\ \mathrm{REPLACES},\\
   \mathrm{EXCLUDES},\ \mathrm{ALTERNATIVE\_TO},\
   \mathrm{MATCHES},\ \mathrm{SUBTRACTS}\}.
\end{gathered}
\label{eq:nuclear-provenance-relations}
\end{equation}
It must reject:
\begin{itemize}
\item an explicit \(NN\pi\) sector together with the same pion contribution
      already contained in the nucleon operator and no subtraction;
\item an explicit compact sector together with the contact operator obtained
      by integrating it out;
\item a coherent double-scattering amplitude and a scalar shadowing ratio
      representing the same contribution;
\item a one-body off-shell correction and the untransformed two-body operator;
\item a wave-function family treated as simultaneous additive states;
\item a CP-map replacement and the amplitude-level mechanism it replaced;
\item partonic soft rescattering and nuclear Glauber exchange in the same
      overlap region without \cref{eq:parton-nuclear-overlap-subtraction}.
\end{itemize}

\subsection{Closure vector}

A compact closure record is
\begin{equation}
 \bm\eps_D=
 (\eps_{\rm norm},\eps_B,\eps_Q,\eps_{P^+},
 \eps_{\rm current},\eps_{\rm ang},
 \eps_{b_1},\eps_{\tagg},\eps_{\rm provenance}).
\label{eq:nuclear-closure-vector}
\end{equation}
Each component has its own tolerance and uncertainty source.  A pass in one
component cannot compensate a failure in another.

\section{Nuclear inference, uncertainty, and falsifiability}
\label{sec:inference}

\subsection{Shared nuclear member}

A nuclear member is
\begin{align}
 \mathfrak m_D=\bigl(&
 \mathfrak m_N,
 H_D,\mathfrak r_D,
 J_D,\widehat\cO_D,
 \Psi_D,
 \text{Fock content},
 \text{coherent order},
 \notag\\
 &\text{matching coefficients},
 \text{subtraction cells},
 \text{numerical realization}
 \bigr).
\label{eq:nuclear-member}
\end{align}
The same identifier propagates through every observable.  A wave-function
choice cannot be varied while silently holding a current or partonic operator
that was fitted only with another choice.

\subsection{Parameter posterior}

Let \(\theta_D\) contain nuclear EFT coefficients, matched partonic-operator
coefficients, coherent-amplitude parameters, and supported compact-sector
couplings.  The posterior is
\begin{equation}
 p(\theta_D,\delta_D\mid D_{\rm cal})
 \propto
 p(D_{\rm cal}\mid\theta_D,\delta_D)
 p(\theta_D)p(\delta_D),
\label{eq:nuclear-posterior}
\end{equation}
where \(\delta_D\) contains separately modeled EFT, Fock, regulator,
matching, and numerical discrepancies.  The calibration set may include:
\begin{itemize}
\item binding energy and asymptotic \(S/D\) information;
\item charge, magnetic, quadrupole, and selected elastic form-factor data;
\item selected inclusive unpolarized and tensor observables;
\item a restricted tagged data subset;
\item diffractive or lattice information needed for coherent/matched
      operators.
\end{itemize}
It must not include one free normalization and transverse width for every
deuteron TMD.

\subsection{Withheld predictions}

Useful withheld families include:
\begin{itemize}
\item tagged tensor observables after calibrating only inclusive quantities;
\item \(LT\) or \(TT\) TMDs after calibrating the \(LL\) tensor sector;
\item deuteron gluon color/helicity channels after quark-sector calibration;
\item form-factor or current moments not used to fit the wave function and
      currents;
\item compact-sector-sensitive observables after fitting only long-distance
      data.
\end{itemize}
A failed withheld family revises the state, Hamiltonian, operator, matching,
or discrepancy model.  It is not repaired by adding a coefficient directly
to the failed TMD.

\subsection{Uncertainty budget}

The reported covariance separates:
\begin{enumerate}
\item nucleon microscopic posterior;
\item nuclear parameter posterior;
\item nuclear EFT truncation;
\item regulator and basis truncation;
\item current and partonic-operator truncation;
\item nuclear Fock-sector truncation;
\item coherent multiple-scattering and diffractive input;
\item parton--nuclear overlap subtraction;
\item TMD/GTMD matching and evolution;
\item numerical integration and transform error.
\end{enumerate}
Alternative high-quality potentials are a controlled model family unless a
common EFT posterior relates them.  Their envelope is not automatically a
confidence interval.

\subsection{Identifiability}

For observables \(\cO_i\) and parameters \(\theta_a\), define
\begin{equation}
 J_{ia}=\frac{\partial\cO_i}{\partial\theta_a}.
\label{eq:nuclear-sensitivity-jacobian}
\end{equation}
Posterior correlations, singular values of \(J\), Fisher information, and
profile likelihoods diagnose whether elastic, inclusive, tagged, and TMD
observables distinguish the proposed mechanisms.  A parameter direction that
only rescales one named TMD has failed the shared-operator objective.

\section{Computational realization and typed interfaces}
\label{sec:software}

\subsection{Core objects}

\begin{longtable}{L{0.27\textwidth}L{0.65\textwidth}}
\toprule
Object & Responsibility\\
\midrule
\endhead
\texttt{NuclearResolution} &
Regulator, basis, EFT order, current/operator order, Fock content, coherent
order, normalization, and scheme.\\
\texttt{NuclearSectorId} &
Typed \(NN\), \(NN\pi\), \(\Delta\Delta\), six-quark, compact, and other
sectors with exact quantum numbers.\\
\texttt{NuclearHamiltonian} &
Hermitian block mass operator, sector-changing interactions, counterterms,
and convergence metadata.\\
\texttt{NuclearFockState} &
Normalized spin-1 sector amplitudes, member identity, phase convention, and
state ledgers.\\
\texttt{LFDeuteronWavefunction} &
Rest-frame interaction, \(S/D\) amplitudes, Jacobian, Melosh/rotation map,
relativistic components, and current family.\\
\texttt{NuclearRecoilMap} &
External and active transfer, intrinsic initial/final coordinates, active
slot, spectator policy, transfer partition, and Jacobian.\\
\texttt{SpectralAmplitude} &
Off-forward deuteron--nucleon helicity matrix and its irreducible and Wigner
reductions.\\
\texttt{MatchedNuclearOperator} &
One-body, two-body, pion, Delta, compact, coherent, and transition blocks with
complete operator identity.\\
\texttt{DeuteronGTMDParent} &
Full sector-summed amplitude-level parent with flavor, helicity, link/color,
rank, transfer, scheme, and member identity.\\
\texttt{CoherentKernel} &
Diffractive helicity amplitudes, propagation phases, scattering order,
Glauber/soft overlap, and factorization status.\\
\texttt{PartialTraceMap} &
Derived CP reduction with explicit amplitude embedding and unresolved
subsystem.\\
\texttt{TaggedObservable} &
Spectator kinematics, spin density, pole variables, FSI status, and
inclusive-sum rule.\\
\texttt{NuclearLedger} &
Normalization, charge, baryon, momentum, current, angular, tensor, tagged,
and provenance residuals.\\
\texttt{NuclearMemberStore} &
Correlated nucleon+nuclear members and covariance-preserving observable
ancestry.\\
\bottomrule
\end{longtable}

\subsection{Typed data flow}

The computational chain is
\begin{equation}
\begin{aligned}
 \texttt{NucleonParent}
 &\longrightarrow \texttt{MatchedNuclearOperator}
 \longrightarrow \texttt{NuclearFockState}\\
 &\longrightarrow \texttt{DeuteronGTMDParent}
 \longrightarrow \texttt{ReductionMap}
 \longrightarrow \texttt{Observable}.
\end{aligned}
\label{eq:nuclear-data-flow}
\end{equation}
Every arrow validates coordinate, sector, operator, scheme, rank, member, and
provenance compatibility.  Array-shape compatibility is insufficient.

\subsection{Immutable cache keys}

A deuteron-parent cache key contains at least
\begin{equation}
 \mathrm{key}=
 \mathrm{hash}(
 \mathfrak m_N,
 \mathfrak m_D,
 \mathfrak o_{a/D},
 \mathfrak r_{\nuc},
 x_D,\kT,\DT,
 \text{reduction},
 \text{code version}).
\label{eq:nuclear-cache-key}
\end{equation}
Changing a wave function, current, matching coefficient, Wilson path, sector,
or recoil convention invalidates the cache.

\subsection{Sparse block structure}

Selection rules determine zero blocks before numerical contraction.  Blocks
are indexed by
\begin{equation}
 (\alpha,\beta,
 \Lambda',\Lambda,
 \lambda_N',\lambda_N,
 J,\text{rank},
 \text{isospin},\text{color},
 \text{link},\text{member}).
\label{eq:nuclear-block-index}
\end{equation}
A block absent by symmetry is not represented by a learned value constrained
to be small.

\subsection{Provenance and ancestry}

Every final value supports the query
\begin{equation}
\begin{aligned}
 \text{observable}
 &\longleftarrow \text{reduction}
 \longleftarrow \text{deuteron parent block}\\
 &\longleftarrow \text{operator block + state sectors}
 \longleftarrow \text{matching and microscopic members}.
\end{aligned}
\label{eq:nuclear-ancestry-query}
\end{equation}
The trace includes excluded alternatives and subtraction partners, not only
active contributions.

\section{Formal requirements and property-based tests}
\label{sec:requirements-tests}

\subsection{Stable requirement identifiers}

\begin{longtable}{L{0.19\textwidth}L{0.71\textwidth}}
\toprule
Identifier & Requirement\\
\midrule
\endhead
\texttt{V4.MATCH.001} & Nuclear Hamiltonian and operators share one matching
and unitary convention.\\
\texttt{V4.MATCH.002} & Eliminated sectors induce both interactions and
operators.\\
\texttt{V4.KIN.001} & \(\DT,\DNT,\pT,\kT,\RT,\bD,\bM\) remain distinct.\\
\texttt{V4.KIN.002} & One-body spectator recoil passes exact physical and
intrinsic closure.\\
\texttt{V4.STATE.001} & All retained nuclear sectors share one normalization
and conserved ledgers.\\
\texttt{V4.STATE.002} & The \(NN\) state records its rest-to-LF map and
current family.\\
\texttt{V4.OP.001} & Every nuclear operator block has complete source/target
sector and matching identity.\\
\texttt{V4.OP.002} & Off-shell one-body and many-body terms transform
covariantly.\\
\texttt{V4.GTMD.001} & The deuteron parent is the full sector matrix element,
not a post hoc sum of named TMDs.\\
\texttt{V4.IA.001} & The IA convolution is derived from the spectator-preserving
one-body block.\\
\texttt{V4.SPIN.001} & Projection occurs after nuclear spin recoupling.\\
\texttt{V4.CURR.001} & Currents satisfy charge, Ward, and angular-condition
tests with the Hamiltonian.\\
\texttt{V4.COH.001} & Polarized/tensor coherent kernels use helicity-resolved
diffractive amplitudes.\\
\texttt{V4.PION.001} & Internal and exchange pion regions are matched and
subtracted once.\\
\texttt{V4.COMP.001} & Compact sectors require explicit amplitudes and
operators before nonzero central use.\\
\texttt{V4.CP.001} & CP maps are derived only after a justified partial trace.\\
\texttt{V4.TAG.001} & Tagged integration recovers inclusive IA and the pole
residue is predicted.\\
\texttt{V4.PROV.001} & Replacement, exclusion, matching, and subtraction
relations are enforceable.\\
\texttt{V4.CONV.001} & Observable convergence is reported along every nuclear
truncation axis.\\
\texttt{V4.PRED.001} & At least one nuclear observable family remains withheld
from calibration.\\
\bottomrule
\end{longtable}

\subsection{Kinematic and state tests}

Mandatory exact or tolerance-controlled tests include:
\begin{enumerate}
\item intrinsic and physical recoil closure;
\item transfer reversal and unit Jacobian;
\item wave-function normalization and exchange antisymmetry;
\item baryon, charge, and plus-momentum ledgers;
\item pure-S, \(D\to0\), and nonrelativistic Melosh limits;
\item spin-1 irreducible reconstruction;
\item Hermiticity of the spectral amplitude;
\item positivity of the forward density without clipping;
\item nuclear mass-operator Hermiticity and eigensolver residuals.
\end{enumerate}

\subsection{Operator and current tests}

\begin{enumerate}
\item operator source/target-sector compatibility;
\item Hermitian-conjugate transition blocks;
\item unitary/field-redefinition invariance of complete matrix elements;
\item current conservation and charge normalization;
\item angular-condition residual and extraction-prescription convergence;
\item GTMD moment versus local-current agreement;
\item one-body/two-body matching-order closure;
\item rejection of incomplete gluon link/color identity.
\end{enumerate}

\subsection{Reduction tests}

\begin{enumerate}
\item full parent to IA reduction in the one-body limit;
\item GTMD-to-TMD, GPD, PDF, form-factor, and current commuting squares;
\item rank-aware Fourier transforms after nuclear composition;
\item \(b_1\) and \(f_{1LL}\) convention adapter;
\item proton/neutron and flavor traceability;
\item tagged-to-inclusive sum rule;
\item zero one-body gluon double-helicity-flip benchmark where applicable;
\item link reversal after nuclear composition.
\end{enumerate}

\subsection{Coherence and provenance tests}

\begin{enumerate}
\item zero diffractive amplitude removes coherent shadowing;
\item unpolarized diffraction cannot silently populate tensor channels;
\item partonic/nuclear soft overlap is subtracted once;
\item \(NN\pi\) internal/exchange pion duplication is rejected;
\item explicit compact sector and integrated-out contact are mutually
      exclusive without a matching cell;
\item coherent interference is not replaced by a CP map;
\item default composition contains exactly one representative of each
      physical mechanism;
\item deliberately duplicated mechanisms trigger a deterministic error
      before numerical evaluation.
\end{enumerate}

\section{Analytic benchmark suite}
\label{sec:benchmarks}

\subsection{Benchmark A: scalar weak-binding two-body state}

Use two distinguishable scalar constituents with normalized wave function
\(\phi(y,\pT)\).  The off-forward density is
\begin{equation}
 \cS(y,\pT,\DT)
 =\phi^*\!\left(y,\pT+\frac{1-y}{2}\DT\right)
 \phi\!\left(y,\pT-\frac{1-y}{2}\DT\right).
\label{eq:benchmark-scalar-density}
\end{equation}
It must satisfy Hermiticity, unit Jacobian, the forward density limit, charge
normalization, and equality of the direct form-factor overlap with the
integrated one-body GTMD.

\subsection{Benchmark B: analytic S--D spin-1 model}

Choose analytic radial functions \(u_0(k)\) and \(u_2(k)\) and construct
\cref{eq:instant-sd-wave}.  The benchmark verifies:
\begin{itemize}
\item exact spin-1 normalization;
\item \(SS,SD,DS,DD\) reconstruction;
\item vanishing nonrelativistic tensor density at \(u_2=0\);
\item linear response of the leading tensor signal to small \(u_2\);
\item correct \(M=0,\pm1,\pm2\) irreducible components.
\end{itemize}
The analytic functions are validation only and cannot enter the central
model as universal transverse profiles.

\subsection{Benchmark C: one-body current and angular condition}

Construct a covariant spin-1 current from three chosen form factors, generate
all plus-current helicity amplitudes, and verify
\cref{eq:angular-condition}.  Then remove one required current component and
confirm a nonzero residual.  This distinguishes a correct extraction
adapter from an incomplete current.

\subsection{Benchmark D: unitary off-shell transformation}

In a two-level nuclear model, apply a unitary transformation
\(U=e^{i\epsilon S}\).  Verify to the declared order that
\begin{equation}
 \langle\Psi|\cO|\Psi\rangle
 =\langle U\Psi|U\cO U^\dagger|U\Psi\rangle,
\label{eq:benchmark-unitary-invariance}
\end{equation}
while the decomposition into one- and two-body terms changes.  Omitting the
induced operator must fail.

\subsection{Benchmark E: \texorpdfstring{\(NN\oplus NN\pi\)}{NN plus NNpi} model}

Use
\begin{equation}
 |D\rangle=\sqrt{Z}|NN\rangle+\sqrt{1-Z}|NN\pi\rangle
\label{eq:benchmark-two-sector-state}
\end{equation}
with a transition operator.  Test total normalization, momentum sharing,
diagonal pion-active contribution, transition interference, and exact
recovery of the integrated-out contact representation after matching.
Enabling both representations must trigger the provenance guard.

\subsection{Benchmark F: coherent double scattering}

Use a two-state diffractive kernel with a known propagation phase.  The
benchmark checks:
\begin{itemize}
\item zero coherent correction when either elementary amplitude vanishes;
\item phase reversal under longitudinal-order interchange;
\item distinct scalar and tensor projections from explicit helicity
      amplitudes;
\item failure of a copied unpolarized shadowing ratio to reproduce the tensor
      result.
\end{itemize}

\subsection{Benchmark G: derived CP map}

Construct an explicit amplitude embedding \(V\) and compare the partial-trace
result with its Kraus representation.  Verify positivity and trace closure.
Then retain a coherent sector superposition and show that tracing before the
amplitudes are combined gives the wrong interference observable.

\subsection{Benchmark H: tagged pole and inclusive recovery}

Use an analytic two-body wave function with a known active-nucleon pole.
Verify:
\begin{enumerate}
\item tagged integration equals inclusive IA;
\item multiplying by the pole factor yields the prescribed residue;
\item a regular FSI model does not alter the pole residue;
\item tensor tagged combinations select the expected D-wave block.
\end{enumerate}

\subsection{Benchmark I: no-double-counting rejection}

Build a composition plan containing:
\begin{enumerate}
\item a nucleon pion-cloud operator;
\item an explicit nuclear pion sector;
\item the required overlap subtraction.
\end{enumerate}
The three-term matched plan passes.  Removing the subtraction or adding a
second scalar pion correction fails deterministically.  Equivalent tests are
performed for compact/contact and coherent/scalar-response pairs.

\section{Contract with Volume V: matching, evolution, and processes}
\label{sec:volume-v-contract}

\subsection{Exported deuteron boundary}

For each parton species, link/color class, target projector, and common member,
Volume~IV exports
\begin{equation}
 \mathfrak W_{a/D}^{(4)}=
 \left(
 W_{a/D,\even}^{[J,\gamma]},
 W_{a/D,\odd}^{[J,\gamma]},
 \mathfrak o_{a/D},
 \mathfrak m_D,
 \mathfrak L_D,
 \mathfrak B_D,
 \operatorname{Cov}_{\micro+\nuc}
 \right),
\label{eq:volume4-export}
\end{equation}
where \(\mathfrak L_D\) is the complete state/current/momentum/provenance
ledger and \(\mathfrak B_D\) includes partonic and nuclear phase/soft overlap
information.  No named scalar TMD is exported without its parent identity and
ancestry.

\subsection{Regulator-to-QCD matching}

Volume~V constructs the declared QCD-scheme object
\begin{equation}
 \widetilde W_{a/D}^{\QCD,[\gamma]}
 (x_D,\bM;\mu,\zeta)
 =\cZ_{a,D}^{\LF\to\TMD}
 (\bM;\mathfrak r_D,\mu,\zeta)
 \otimes
 \widetilde W_{a/D}^{(4),[\gamma]}
 +\delta W_{a/D}^{\rm power}.
\label{eq:deuteron-lf-to-tmd-matching}
\end{equation}
The matching coefficient includes any operator mixing and identifies which
nuclear and nucleon regulator dependences cancel.  It does not replace the
nuclear matrix element by an external deuteron fit.

\subsection{When evolution commutes with nuclear composition}

At leading power, the Collins--Soper evolution kernel depends on the parton
color representation and not on target polarization.  Let
\(\cU_a(\bM;\mu,\zeta;\mu_0,\zeta_0)\) be a linear multiplicative evolution
operator in \(\bM\) space.  For a one-body nuclear map \(\cN_{\IA}\) that
contains no additional \(\mu,\zeta\) dependence,
\begin{equation}
 \cU_a\,\cN_{\IA}[\widetilde W_{a/N}]
 =\cN_{\IA}[\cU_a\widetilde W_{a/N}].
\label{eq:cs-evolution-commutes-ia}
\end{equation}
This equality is conditional on a common scheme, scales, link class, soft
subtraction, and parton representation.  It need not hold term by term for
operator mixing, two-body matching coefficients, coherent kernels, or
subtractions.

\begin{requirement}[Evolution commuting square]
\label{req:evolution-commuting-square}
For every nuclear operator class, Volume~V must either demonstrate the
commuting square between nuclear composition and evolution/matching or retain
an explicit noncommuting correction.  The target independence of the
Collins--Soper kernel is not permission to erase nuclear operator mixing.
\end{requirement}

\subsection{Collinear evolution and hadronic sectors}

The full collinear boundary includes quark, antiquark, and gluon content in
all nuclear sectors.  Nonsinglet, singlet, and quark--gluon evolution act on
the matched partonic vector.  Pion or Delta partons are not evolved as an
additional copy after they have already been matched into the deuteron
quark/gluon distributions.  The hadronic-sector ledger remains provenance;
the evolved QCD object is the sum defined by the matched operator.

\subsection{Process assembly}

A process object consumes \cref{eq:volume4-export} only if its factorization
status permits the requested link/color structure.  The observable has the
schematic form
\begin{equation}
 \frac{\dd\sigma}{\dd\mathcal P}
 =W_{\rm TMD}(\qT,Q)
 +Y(\qT,Q),
 \qquad
 Y=\sigma_{\rm FO}-[W_{\rm TMD}]_{\rm expanded\ to\ FO},
\label{eq:volume5-wy-contract}
\end{equation}
with explicit hard, fragmentation, color, and target-polarization projectors.
Nuclear high-momentum structure is not tuned to reproduce the fixed-order
large-\(q_T\) tail.

\subsection{Resolved evolution and closure}

Evolution must preserve the resolved decomposition
\begin{equation}
 \{
 p\text{-in-}D,
 n\text{-in-}D,
 p+n,
 p-n,
 \text{many-body correction},
 \text{total}
 \}
\label{eq:resolved-evolution-components}
\end{equation}
for every flavor, target channel, rank, link/color class, mechanism, and
correlated member.  Closure after evolution is tested against the evolved
total, not assumed from the unevolved boundary.

\section{Staged C4 implementation program}
\label{sec:c4-program}

\subsection{Relation to the C3 pilot}

C3 constructs a validation-only zero-skewness nucleon overlap pilot and keeps
it disconnected from the accepted central artifacts until its analytic gates
pass.  C4 must consume only a C3 object that satisfies the complete Volume~II
contract.  It must not connect the accepted phenomenological nucleon tables to
a new nuclear state and call the result microscopic.

The exact repository-specific C4 prompt should be written after the C3 report
is available.  The formal work packages below define the intended order.

\subsection{C4A: nuclear coordinates and two-body recoil}

Implement:
\begin{itemize}
\item deuteron incoming/outgoing momentum fibers;
\item \(x_D\leftrightarrow x_N\) adapter;
\item two-body intrinsic configurations and measures;
\item spectator-preserving recoil from
      \cref{eq:active-intrinsic-recoil,eq:spectator-intrinsic-recoil};
\item transfer-partition objects for two-body operators;
\item exact closure, inverse, Jacobian, and mismatch tests.
\end{itemize}
No wave function or GTMD is introduced until this layer passes.

\subsection{C4B: spin-1 light-front wave function and spectral amplitude}

Implement a validation-only \(S/D\) state, followed by one declared realistic
wave-function adapter.  Required outputs are:
\begin{itemize}
\item normalized helicity amplitudes;
\item rest-to-light-front Jacobian and spin rotations;
\item irreducible \(K=0,1,2\) density components;
\item off-forward spectral amplitude and nuclear Wigner transform;
\item \(SS,SD,DS,DD\) ancestry;
\item pure-S, nonrelativistic, Hermiticity, and positivity tests.
\end{itemize}
The central accepted boundary remains disconnected.

\subsection{C4C: impulse common-parent pilot}

Compose the C3 nucleon parent with the C4B spectral amplitude using the native
typed reduction/provenance infrastructure.  Validate:
\begin{itemize}
\item \cref{eq:ia-gtmd-convolution} at nonzero \(\DT\);
\item TMD, GPD, PDF, and current reductions;
\item proton/neutron and flavor traceability;
\item spin recoupling before projection;
\item rank-aware Fourier behavior;
\item exact isolation from the accepted production root.
\end{itemize}

\subsection{C4D: local currents, form factors, \texorpdfstring{\(b_1\)}{b1}, and tagging}

Add the minimal consistent current and projector layer required to test:
\begin{itemize}
\item charge, magnetic, and quadrupole normalization;
\item the angular condition;
\item GTMD moment/current agreement;
\item the tensor PDF and \(b_1\) relation;
\item tagged-to-inclusive closure and an analytic pole benchmark.
\end{itemize}
This package is still a pilot and does not fit new nuclear coefficients.

\subsection{C4E: matched two-sector and pion pilot}

Introduce an analytic \(NN\oplus NN\pi\) model with one transition operator.
Demonstrate:
\begin{itemize}
\item common normalization and plus-momentum closure;
\item pion-active and nucleon-active operator blocks;
\item transition interference;
\item explicit-pion versus integrated-out-contact equivalence;
\item deterministic rejection of double counting.
\end{itemize}

\subsection{C4F: coherent amplitude and CP-reduction pilot}

Implement a two-step helicity-resolved coherent kernel and a separately
derived partial-trace CP map.  Prove that:
\begin{itemize}
\item the coherent tensor result is not a copied unpolarized ratio;
\item zero elementary diffraction removes the correction;
\item partonic/nuclear soft overlap is represented explicitly;
\item tracing before coherent amplitudes are combined fails the benchmark;
\item the final provenance and closure manifests are deterministic.
\end{itemize}

\subsection{Production isolation}

Until all C4 analytic and current gates pass:
\begin{enumerate}
\item no C4 object is reachable from the accepted production composition
      root;
\item the 216 accepted C2 reductions remain unchanged;
\item authoritative parent/correlator artifacts remain byte-identical;
\item pilot parameters are labeled \texttt{VALIDATION\_ONLY};
\item the pilot cannot be selected by a production configuration.
\end{enumerate}
A later work package may authorize a physics comparison, but not silent
replacement.

\section{Acceptance criteria and scientific status}
\label{sec:acceptance}

\subsection{Volume IV formal acceptance}

The formalism is internally complete only when:
\begin{enumerate}
\item The quark--gluon-to-nuclear hierarchy, model space, regulator, and
      matching order are declared.
\item Nuclear Hamiltonian, currents, and partonic operators transform as one
      matched family.
\item The spin-1 state contains normalized \(NN\) and every supported
      additional sector with one charge and momentum ledger.
\item The \(NN\) light-front wave function has an explicit Jacobian,
      spin-rotation convention, relativistic-component status, and current
      family.
\item The one-body recoil map is exact and two-body transfer sharing is
      operator defined.
\item The off-forward spectral amplitude is a full deuteron--nucleon helicity
      matrix and satisfies Hermiticity and forward positivity.
\item The deuteron parent is evaluated as
      \cref{eq:full-deuteron-parent} with sector-diagonal and allowed
      transition blocks.
\item Impulse approximation is derived from the full parent and its TMD, GPD,
      PDF, form-factor, and current reductions commute within stated errors.
\item Spin-1 projections occur after nucleon spin recoupling, with exact
      convention adapters.
\item Elastic current normalization and the angular condition pass with
      currents consistent with the Hamiltonian.
\item \(b_1\), tensor PDFs, and conditional tensor sum rules carry their
      exact assumptions.
\item Coherent shadowing uses helicity-resolved amplitudes; tensor channels are
      not copied from unpolarized diffraction.
\item The partonic Wilson/soft sector and nuclear coherent sector have an
      explicit overlap subtraction when required.
\item The \(NN\pi\) sector has a spin-resolved amplitude, active-pion and
      transition operators, and a pion double-counting subtraction.
\item \(\Delta\Delta\), compact, and hidden-color sectors remain zero-centered
      unless their full amplitude/operator requirements are met.
\item CP maps are derived by justified partial traces and never replace
      unresolved coherence.
\item Tagged integration closes to inclusive IA and the pole residue is
      predicted by the same spectral amplitude.
\item State, current, partonic, tensor, tagged, and provenance ledgers pass
      separately.
\item Observable convergence is demonstrated along all retained nuclear
      truncation axes.
\item At least one nontrivial deuteron observable family is withheld from
      calibration and predicted by the shared member ensemble.
\end{enumerate}

\subsection{Predictive status}

Completion of Volume~IV supplies a matched spin-1 nuclear parent but does not
alone establish a complete prediction.  The model becomes genuinely
predictive only after:
\begin{itemize}
\item the nucleon common state and Wilson dynamics have passed Volumes~I--III;
\item the nuclear hierarchy passes the acceptance criteria above;
\item Volume~V matches the full parent to one QCD scheme, evolves all ranks and
      partonic sectors, and constructs valid process kernels;
\item correlated calibration and withheld validation succeed without
      TMD-specific repair parameters.
\end{itemize}

\subsection{Recommended foundation}

The adopted nuclear foundation is
\begin{quote}
\emph{a regulated, spin-resolved nuclear light-front EFT matched to a common
quark--gluon nucleon theory, with one normalized deuteron Fock state,
Hamiltonian-consistent currents and partonic operators, exact typed recoil,
amplitude-level coherent and sector interference, derived reduced channels,
and observable-level convergence and falsifiability.}
\end{quote}

\section{Conclusion}
\label{sec:conclusion}

The central result of this volume is that nuclear composition is not a final
smearing step.  It is a second microscopic layer with its own Hilbert space,
Hamiltonian, currents, operator matching, coherence, and truncation theory.
The deuteron GTMD is therefore not a nucleon TMD multiplied by a tensor factor
or supplemented by independently normalized nuclear corrections.  It is the
matrix element of a matched quark/gluon operator in one normalized spin-1
state.

The weak-binding \(NN\) component supplies a controlled and highly constrained
starting point.  Its exact light-front recoil map, \(S/D\) interference,
Melosh rotations, spin-density matrix, elastic currents, tagged observables,
and impulse GTMD reductions form a stringent validation system.  The same
formalism shows where the impulse approximation is insufficient: coherent
small-\(x\) scattering, two-body currents, pion and Delta sectors,
short-distance contacts, and compact color structure require new operator
blocks or explicit state components.  Their contributions are correlated by
matching and cannot be introduced as unrelated functions.

Volume~IV also resolves several recurrent ambiguities.  Active-nucleon
transfer is distinct from the intrinsic deuteron wave-function shift.
Off-shell one-body corrections are representation dependent and meaningful
only with their compensating many-body operators.  Completely positive maps
are reduced descriptions derived after a trace, not substitutes for coherent
amplitudes.  Partonic Wilson transport and nuclear rescattering are separate
objects whose overlap must be subtracted.  Pion and compact sectors may be
represented explicitly or through induced operators, but not both without a
matching cell.

The next formal layer is Volume~V: common regulator-to-QCD matching,
rank-aware evolution, collinear mixing, and process-level \(W+Y\) assembly.
The immediate computational task, after the C3 report, is a repository-specific
C4 prompt implementing the isolated nuclear recoil, spin-1 wave-function,
spectral-amplitude, and impulse-parent pilot defined here.

\appendix

\section{Derivation of the spectator-preserving recoil map}
\label{app:recoil-derivation}

Let a two-body state of total transverse momentum \(P_T\) have physical
constituent momenta
\begin{equation}
 p_{1T}=yP_T+\pT,
 \qquad
 p_{2T}=(1-y)P_T-\pT.
\label{eq:two-body-physical-momenta}
\end{equation}
In the symmetric frame,
\(
 P_{iT}=-\DT/2
\) and
\(
 P_{fT}=+\DT/2
\).
Demand that the spectator physical momentum be unchanged:
\begin{equation}
 p_{2T}^{\rm in}=p_{2T}^{\rm out}=-\pT.
\label{eq:spectator-fixed-condition}
\end{equation}
Solving
\begin{align}
 -\pT&=(1-y)\left(-\frac\DT2\right)+\bm p_{2T}^{\rm in},
\notag\\
 -\pT&=(1-y)\left(+\frac\DT2\right)+\bm p_{2T}^{\rm out}
\end{align}
gives \cref{eq:spectator-intrinsic-recoil}.  Intrinsic closure then fixes the
active arguments in \cref{eq:active-intrinsic-recoil}.  The active physical
momenta follow as
\begin{align}
 p_{1T}^{\rm in}
 &=y\left(-\frac\DT2\right)
 +\pT-\frac{1-y}{2}\DT
 =\pT-\frac\DT2,
\notag\\
 p_{1T}^{\rm out}
 &=y\left(+\frac\DT2\right)
 +\pT+\frac{1-y}{2}\DT
 =\pT+\frac\DT2.
\end{align}
Thus \(p_{1T}^{\rm out}-p_{1T}^{\rm in}=\DT\).  The transformation is a
translation of the intrinsic variables and has unit Jacobian.

\section{Spin-1 helicity and tensor adapters}
\label{app:spin-adapters}

For helicity-diagonal quantities define
\begin{align}
 F_U&=\frac13(F_++F_0+F_-),
\label{eq:FU-adapter}\\
 F_L&=\frac12(F_+-F_-),
\label{eq:FL-adapter}\\
 \delta_TF&=F_0-\frac12(F_++F_-).
\label{eq:deltaT-adapter}
\end{align}
The inverse is
\begin{align}
 F_0&=F_U+\frac23\delta_TF,
\label{eq:F0-inverse}\\
 F_+&=F_U+F_L-\frac13\delta_TF,
\label{eq:Fplus-inverse}\\
 F_-&=F_U-F_L-\frac13\delta_TF.
\label{eq:Fminus-inverse}
\end{align}
These relations are independent of the named \(S_{LL}\) convention.  The
project adapter \(f_{1LL}=-(2/3)\delta_Tf_1\) is applied afterward.

For a complete matrix, use the irreducible basis in
\cref{eq:spin1-irrep-basis}.  A real Cartesian basis is obtained by
\begin{align}
 T_{K,M}^{c}&=\frac{T_M^{(K)}+(-1)^M T_{-M}^{(K)}}{\sqrt2},
\notag\\
 T_{K,M}^{s}&=\frac{T_M^{(K)}-(-1)^M T_{-M}^{(K)}}{i\sqrt2},
 \qquad M>0,
\label{eq:real-spherical-adapter}
\end{align}
plus \(T_0^{(K)}\).  The mapping to \(L,T_x,T_y,LL,LT_x,LT_y,TT_x,TT_y\)
is stored as an invertible matrix and tested by Gram reconstruction.

\section{Unitary redistribution of one- and two-body operators}
\label{app:unitary-redistribution}

Let
\begin{equation}
 U=e^{i\epsilon S},
 \qquad S=S^\dagger.
\label{eq:unitary-small-transform}
\end{equation}
To first order,
\begin{align}
 H'&=H+i\epsilon[S,H]+\cO(\epsilon^2),
\notag\\
 \cO'&=\cO+i\epsilon[S,\cO]+\cO(\epsilon^2),
\notag\\
 |\Psi'\rangle&=|\Psi\rangle+i\epsilon S|\Psi\rangle+\cO(\epsilon^2).
\label{eq:unitary-first-order}
\end{align}
Then
\begin{align}
 \langle\Psi'|\cO'|\Psi'\rangle
 ={}&\langle\Psi|\cO|\Psi\rangle
 +i\epsilon\langle\Psi|[S,\cO]|\Psi\rangle
 \notag\\
 &-i\epsilon\langle\Psi|S\cO|\Psi\rangle
 +i\epsilon\langle\Psi|\cO S|\Psi\rangle
 +\cO(\epsilon^2)
\notag\\
 ={}&\langle\Psi|\cO|\Psi\rangle+\cO(\epsilon^2).
\label{eq:unitary-cancellation}
\end{align}
If \(S\) mixes one- and two-body sectors, the commutator changes the apparent
body order of \(\cO\).  Keeping the transformed state but dropping the
induced term violates \cref{eq:unitary-cancellation} and creates a spurious
off-shell dependence.

\section{Machine-readable schemas}
\label{app:schemas}

\subsection{Nuclear member record}

\begin{verbatim}
{
  "schema": "deuteron_wigner.volume4.nuclear_member.v1",
  "nucleon_member_id": "...",
  "nuclear_hamiltonian_id": "...",
  "nuclear_resolution_id": "...",
  "current_family_id": "...",
  "partonic_operator_family_id": "...",
  "fock_sector_ids": ["NN", "NNPI", "DELTA_DELTA", "SIX_QUARK"],
  "matching_record_ids": ["..."],
  "subtraction_cell_ids": ["..."],
  "coherent_kernel_id": "...",
  "numerical_realization_id": "...",
  "provenance_node": "..."
}
\end{verbatim}

\subsection{Nuclear recoil record}

\begin{verbatim}
{
  "schema": "deuteron_wigner.volume4.nuclear_recoil.v1",
  "incoming_deuteron_fiber": "...",
  "outgoing_deuteron_fiber": "...",
  "external_delta_t": [0.0, 0.0],
  "active_nucleon_delta_t": [0.0, 0.0],
  "transfer_partition_t": null,
  "y": 0.5,
  "p_t": [0.0, 0.0],
  "active_slot": 0,
  "spectator_policy": "PRESERVED",
  "intrinsic_in": ["..."],
  "intrinsic_out": ["..."],
  "jacobian_id": "...",
  "convention_version": "V4.KIN.002"
}
\end{verbatim}

\subsection{Matched nuclear operator record}

\begin{verbatim}
{
  "schema": "deuteron_wigner.volume4.matched_nuclear_operator.v1",
  "operator_id": "...",
  "parton_species": "QUARK|ANTIQUARK|GLUON",
  "projection_id": "...",
  "source_sector": "NN",
  "target_sector": "NN",
  "body_order": "ONE_BODY|TWO_BODY|PION|DELTA|SHORT|COHERENT|TRANSITION",
  "active_constituent": "PROTON|NEUTRON|PION|DELTA|COMPACT",
  "wilson_path_id": "...",
  "gluon_link_pair_id": "...",
  "gluon_color_class": "F_TYPE|D_TYPE|NOT_APPLICABLE",
  "nuclear_recoil_id": "...",
  "matching_scale": "...",
  "regulator_id": "...",
  "scheme_id": "...",
  "member_id": "...",
  "subtraction_partner_ids": ["..."],
  "provenance_node": "..."
}
\end{verbatim}

\subsection{Deuteron parent record}

\begin{verbatim}
{
  "schema": "deuteron_wigner.volume4.deuteron_gtmd_parent.v1",
  "operator_id": "...",
  "nuclear_member_id": "...",
  "x_coordinate": "X_D",
  "x": 0.1,
  "k_t": [0.0, 0.0],
  "delta_t": [0.0, 0.0],
  "target_helicity_matrix": "...",
  "parton_helicity_matrix": "...",
  "sector_block_contributions": ["..."],
  "reduction_capabilities": ["TMD", "WIGNER", "GPD", "PDF", "CURRENT"],
  "closure_ledger_id": "...",
  "provenance_root": "...",
  "status": "VALIDATION_ONLY|AUTHORITATIVE"
}
\end{verbatim}

\subsection{Nuclear closure record}

\begin{verbatim}
{
  "schema": "deuteron_wigner.volume4.nuclear_closure.v1",
  "member_id": "...",
  "normalization_residual": 0.0,
  "baryon_residual": 0.0,
  "charge_residual": 0.0,
  "hadronic_momentum_residual": 0.0,
  "partonic_momentum_residual": 0.0,
  "current_residual": 0.0,
  "angular_condition_residual": 0.0,
  "b1_residual": 0.0,
  "tagged_inclusive_residual": 0.0,
  "provenance_residual": 0.0,
  "tolerances": {"...": 0.0},
  "status": "PASS|FAIL"
}
\end{verbatim}

\section{Compact acceptance checklist}
\label{app:checklist}

Before a Volume~IV result is used in matching, evolution, or inference, answer:
\begin{enumerate}
\item Are the nuclear Hamiltonian, currents, and partonic operators one matched
      family?
\item Is the deuteron state normalized across all retained sectors?
\item Are external transfer, active-nucleon transfer, intrinsic recoil, and
      partonic momentum distinct and validated?
\item Is the rest-to-light-front wave-function map explicit?
\item Does the spectral amplitude retain deuteron and nucleon helicities?
\item Is every operator block typed by source/target sector and body order?
\item Is the full deuteron parent evaluated before mechanism or named-TMD
      projection?
\item Do TMD, GPD, PDF, current, and form-factor reductions commute within
      reported errors?
\item Do charge, momentum, elastic, angular, tensor, and tagged ledgers pass?
\item Are off-shell and many-body terms representation covariant?
\item Are coherent tensor amplitudes calculated rather than copied from an
      unpolarized ratio?
\item Are partonic and nuclear soft/Glauber overlaps subtracted once?
\item Is the \(NN\pi\) sector spin resolved and pion double counting excluded?
\item Do compact sectors have complete amplitudes and operators before nonzero
      central use?
\item Are CP maps derived only after a justified trace?
\item Does the provenance graph reject every duplicate representation?
\item Are all nuclear truncation axes reported separately?
\item Is at least one deuteron observable family genuinely withheld?
\end{enumerate}
A negative answer blocks authoritative use and identifies the next physics,
matching, or computational task.

\small
\begin{thebibliography}{99}

\bibitem{Dirac1949}
P.~A.~M.~Dirac,
``Forms of relativistic dynamics,''
Rev.\ Mod.\ Phys.\ \textbf{21}, 392 (1949).

\bibitem{KeisterPolyzou1991}
B.~D.~Keister and W.~N.~Polyzou,
``Relativistic Hamiltonian dynamics in nuclear and particle physics,''
Adv.\ Nucl.\ Phys.\ \textbf{20}, 225 (1991).

\bibitem{Carbonell1998}
J.~Carbonell, B.~Desplanques, V.~A.~Karmanov, and J.-F.~Mathiot,
``Explicitly covariant light-front dynamics and relativistic few-body
systems,''
Phys.\ Rept.\ \textbf{300}, 215 (1998), arXiv:nucl-th/9804029.

\bibitem{ChungCoesterKeisterPolyzou1988}
P.~L.~Chung, F.~Coester, B.~D.~Keister, and W.~N.~Polyzou,
``Hamiltonian light-front dynamics of elastic electron--deuteron scattering,''
Phys.\ Rev.\ C \textbf{37}, 2000 (1988).

\bibitem{CookeMiller2002}
J.~R.~Cooke and G.~A.~Miller,
``Deuteron binding energies and form factors from light-front Hamiltonian
field theory,''
Phys.\ Rev.\ C \textbf{66}, 034002 (2002), arXiv:nucl-th/0112037.

\bibitem{Wiringa1995}
R.~B.~Wiringa, V.~G.~J.~Stoks, and R.~Schiavilla,
``Accurate nucleon--nucleon potential with charge-independence breaking,''
Phys.\ Rev.\ C \textbf{51}, 38 (1995), arXiv:nucl-th/9408016.

\bibitem{Machleidt2001}
R.~Machleidt,
``The high-precision, charge-dependent Bonn nucleon--nucleon potential,''
Phys.\ Rev.\ C \textbf{63}, 024001 (2001), arXiv:nucl-th/0006014.

\bibitem{Piarulli2016}
M.~Piarulli et al.,
``Local chiral potentials and the structure of light nuclei,''
Phys.\ Rev.\ C \textbf{94}, 054007 (2016), arXiv:1606.06335.

\bibitem{Krebs2020}
H.~Krebs,
``Nuclear currents in chiral effective field theory,''
Eur.\ Phys.\ J.\ A \textbf{56}, 234 (2020), arXiv:2008.00974.

\bibitem{Schiavilla2018}
R.~Schiavilla et al.,
``Local chiral interactions and magnetic structure of few-nucleon systems,''
Phys.\ Rev.\ C \textbf{99}, 034005 (2019), arXiv:1809.10180.

\bibitem{FurnstahlHammerTirfessa2001}
R.~J.~Furnstahl, H.-W.~Hammer, and N.~Tirfessa,
``Field redefinitions at finite density,''
Nucl.\ Phys.\ A \textbf{689}, 846 (2001), arXiv:nucl-th/0010078.

\bibitem{CanoPire2002}
F.~Cano and B.~Pire,
``Deeply virtual Compton scattering on spin-1 nuclei,''
arXiv:hep-ph/0211444.

\bibitem{CosynPire2018}
W.~Cosyn and B.~Pire,
``Transversity generalized parton distributions for the deuteron,''
Phys.\ Rev.\ D \textbf{98}, 074020 (2018), arXiv:1806.01177.

\bibitem{CosynPire2019}
W.~Cosyn and B.~Pire,
``Deuteron helicity-flip generalized parton distributions in a convolution
model,''
arXiv:1902.03681.

\bibitem{HoodbhoyJaffeManohar1989}
P.~Hoodbhoy, R.~L.~Jaffe, and A.~Manohar,
``Novel effects in deep-inelastic scattering from spin-one hadrons,''
Nucl.\ Phys.\ B \textbf{312}, 571 (1989).

\bibitem{HERMES2005}
A.~Airapetian et al.\ (HERMES Collaboration),
``First measurement of the tensor structure function \(b_1\) of the
deuteron,''
Phys.\ Rev.\ Lett.\ \textbf{95}, 242001 (2005), arXiv:hep-ex/0506018.

\bibitem{CosynDongKumanoSargsian2017}
W.~Cosyn, Y.-B.~Dong, S.~Kumano, and M.~Sargsian,
``Tensor-polarized structure function \(b_1\) in the standard convolution
description of the deuteron,''
Phys.\ Rev.\ D \textbf{95}, 074036 (2017), arXiv:1702.05337.

\bibitem{Kumano2014}
S.~Kumano,
``Tensor-polarized structure functions: Tensor structure of deuteron in the
2020s,''
J.\ Phys.\ Conf.\ Ser.\ \textbf{543}, 012001 (2014), arXiv:1407.3852.

\bibitem{KumanoSong2021}
S.~Kumano and Q.-T.~Song,
``Twist-2 relation and sum rule for tensor-polarized parton distribution
functions of spin-1 hadrons,''
JHEP \textbf{09}, 141 (2021), arXiv:2106.15849.

\bibitem{FrankfurtGuzeyStrikman2012}
L.~Frankfurt, V.~Guzey, and M.~Strikman,
``Leading twist nuclear shadowing phenomena in hard processes with nuclei,''
Phys.\ Rept.\ \textbf{512}, 255 (2012), arXiv:1106.2091.

\bibitem{Miller2014}
G.~A.~Miller,
``Pionic and hidden-color, six-quark contributions to the deuteron \(b_1\)
structure function,''
Phys.\ Rev.\ C \textbf{89}, 045203 (2014), arXiv:1311.4561.

\bibitem{MillerHiddenColor2014}
G.~A.~Miller,
``Hidden color and the \(b_1\) structure function of the deuteron,''
arXiv:1412.1864.

\bibitem{NeffFeldmeier2016}
T.~Neff and H.~Feldmeier,
``The Wigner function and short-range correlations in the deuteron,''
arXiv:1610.04066.

\bibitem{CosynSargsianWeiss2016}
W.~Cosyn, M.~Sargsian, and C.~Weiss,
``Tagged spectator DIS on a polarized spin-1 target,''
PoS \textbf{DIS2016}, 210 (2016), arXiv:1609.01970.

\bibitem{CosynWeiss2020}
W.~Cosyn and C.~Weiss,
``Polarized electron--deuteron deep-inelastic scattering with spectator
nucleon tagging,''
Phys.\ Rev.\ C \textbf{102}, 065204 (2020), arXiv:2006.03033.

\bibitem{CosynWeiss2026I}
W.~Cosyn and C.~Weiss,
``Semi-inclusive deep-inelastic scattering on a polarized spin-1 target. I.
Cross section and spin observables,''
arXiv:2603.23699 (2026).

\bibitem{CosynWeiss2026II}
W.~Cosyn and C.~Weiss,
``Semi-inclusive deep-inelastic scattering on a polarized spin-1 target. II.
Deuteron and spectator nucleon tagging,''
arXiv:2603.23700 (2026).

\bibitem{CaoEtAl2026}
X.~Cao, S.~Cheng, Y.~Duan, Y.~Li, S.~Xu, and X.~Zhao,
``Non-perturbative flavor asymmetry in the nucleon and deuteron: The
light-front Hamiltonian effective field theory approach,''
arXiv:2601.13567 (2026).

\bibitem{BacchettaMulders2000}
A.~Bacchetta and P.~J.~Mulders,
``Deep inelastic leptoproduction of spin-one hadrons,''
Phys.\ Rev.\ D \textbf{62}, 114004 (2000), arXiv:hep-ph/0007120.

\bibitem{KumanoSongTMD2021}
S.~Kumano and Q.-T.~Song,
``Transverse-momentum-dependent parton distribution functions for spin-1
hadrons,''
Phys.\ Rev.\ D \textbf{103}, 014025 (2021), arXiv:2011.08583.

\bibitem{Boer2016}
D.~Boer, S.~Cotogno, T.~van Daal, P.~J.~Mulders, A.~Signori, and Y.-J.~Zhou,
``Gluon and Wilson-loop TMDs for hadrons of spin \(\leq1\),''
JHEP \textbf{10}, 013 (2016), arXiv:1607.01654.

\bibitem{Cotogno2017}
S.~Cotogno, T.~van Daal, and P.~J.~Mulders,
``Positivity bounds on gluon TMDs for hadrons of spin \(\leq1\),''
JHEP \textbf{11}, 185 (2017), arXiv:1709.07827.

\bibitem{Collins2011}
J.~C.~Collins,
\emph{Foundations of Perturbative QCD},
Cambridge University Press (2011).

\bibitem{EchevarriaGTMD2016}
M.~G.~Echevarria, A.~Idilbi, K.~Kanazawa, C.~Lorc\'e, A.~Metz,
B.~Pasquini, and M.~Schlegel,
``Proper definition and evolution of generalized transverse momentum
distributions,''
Phys.\ Lett.\ B \textbf{759}, 336 (2016), arXiv:1602.06953.

\end{thebibliography}

\end{document}
